\documentclass[11pt,a4paper]{article}

\usepackage[utf8]{inputenc}
\usepackage[T1]{fontenc}
\usepackage{geometry}
\geometry{margin=1in}

% Paragraph spacing
\setlength{\parskip}{0.6\baselineskip plus 2pt minus 1pt}
\setlength{\parindent}{0pt}

\usepackage{mathtools,amsmath,amssymb,amsthm}
\usepackage{booktabs}
\usepackage[hidelinks]{hyperref}
\usepackage{enumitem}
\usepackage{url}
\usepackage{cleveref}

% Theorem environments
\theoremstyle{plain}
\newtheorem{theorem}{Theorem}[section]
\newtheorem{lemma}[theorem]{Lemma}
\newtheorem{corollary}[theorem]{Corollary}
\newtheorem{proposition}[theorem]{Proposition}
\theoremstyle{definition}
\newtheorem{definition}[theorem]{Definition}
\newtheorem{axiom}[theorem]{Axiom}
\newtheorem{assumption}[theorem]{Assumption}
\newtheorem{example}[theorem]{Example}
\theoremstyle{remark}
\newtheorem{remark}[theorem]{Remark}

\newcommand{\defeq}{\vcentcolon=}
\newcommand{\nf}{\mathrm{nf}}
% Coq theorem names (typewriter, breakable at underscores)
\DeclareUrlCommand\coq{\urlstyle{tt}}
\Urlmuskip=0mu
\emergencystretch=3em

\title{Normalization Confluence for\\Registry-Governed Stream Processing}

\author{
Dayna Blackwell\\
\small\href{mailto:dayna@blackwell-systems.com}{dayna@blackwell-systems.com}\\
\small Licensed under CC-BY-4.0.}

\date{4 October 2026\\[4pt]
\small Version 2. See Appendix (Errata and corrections) for changes from version 1.}

\begin{document}
\maketitle

\begin{abstract}
\setlength{\parskip}{0.6\baselineskip plus 2pt minus 1pt}
Coordination-free convergence in distributed systems has been achieved through
two structural regimes: \emph{operation commutativity} (CRDTs), where
operations commute by construction, and \emph{invariant confluence}
(I-confluence), where all orderings preserve application invariants.  We
identify a third regime---\emph{normalization confluence}---in which
operations may be non-commutative and may violate invariants, but convergence
is restored through a compensation operator that normalizes states to
validity.

We formalize this regime by introducing \emph{registry-governed stream
processing}, where an external registry of invariant predicates defines valid
states and a compensation operator repairs violations.  We construct a
\emph{governance rewrite system} whose configurations pair a state with a
buffer of pending events, and whose rewrite rules include both event
application (in any causality-respecting order) and compensation.  Two
structural conditions---\emph{well-founded compensation} (WFC) and
\emph{compensation commutativity} (CC)---ensure termination and confluence
respectively.  Under both, Newman's Lemma guarantees unique normal forms: any
two processors consuming the same events agree on a valid state, regardless
of the order in which they apply those events.  A third condition---%
\emph{uniformly bounded compensation} (UBC)---yields constant per-event
overhead and $O(n \log n)$ total cost matching conventional stream processing.
We show UBC is tight: without it, compensation depth grows without limit.
We show CC is exact where it can be: with every buffered event free to fire
and compensation that repairs to validity, every event set delivered from the
initial state has a unique normal form if and only if CC holds on the
reachable states; when causal dependencies constrain the future of a run, the
exact condition is a joinability form of CC.

Compensation commutativity is strictly weaker than operation commutativity:
operations need not commute, only their compensated results.  The gap between
these two conditions is precisely the design space that normalization
confluence occupies.

To make CC practically verifiable, we develop a \emph{verification calculus}
that reduces the global CC obligation to local checks.  An \emph{invariant
footprint} mapping identifies which invariants each event can affect, and a
\emph{decomposable repair} condition ensures per-invariant repairs commute.
Under decomposable repair, CC2 is equivalent to a per-event strong absorption
property, and CC1 for a pair reduces to commutation of the two events up to
normalization plus a per-event footprint absorption check; at valid states,
disjoint footprints and repair locality reduce CC1 to commutation up to
normalization alone.  A product composition theorem enables modular
verification of multi-registry systems.  The convergence theorems, their
exact converses, the calculus, and the paper's examples and counterexamples
are mechanized in Coq/Rocq, axiom-free.
\end{abstract}

\medskip\noindent\textbf{Keywords:} Stream Processing, Normalization
Confluence, Semantic Consistency, Rewrite Systems, Newman's Lemma,
Compensation Commutativity, Distributed Systems

\clearpage

%========================================
\section{Introduction}
\label{sec:intro}

Coordination-free convergence---the ability of distributed processors to agree
on a result without explicit synchronization---is among the central goals of
distributed systems design.  The literature has identified two structural
regimes under which it is achievable:

\begin{enumerate}[leftmargin=*,nosep]
\item \textbf{Operation commutativity} (CRDTs~\cite{shapiro2011crdt}):
  Operations are designed to commute algebraically.  Any application order
  produces the same state.
\item \textbf{Invariant confluence}
  (I-confluence~\cite{bailis2014coordination}): Operations individually
  preserve application invariants under all orderings.  No repair is needed
  because validity is never violated.
\end{enumerate}

Both regimes achieve convergence by \emph{preventing} inconsistency.  CRDTs
prevent it through algebraic design; I-confluence prevents it through
invariant-preserving operation selection.  Neither accommodates operations that
are non-commutative \emph{and} may violate invariants---yet such operations
are pervasive in practice.  An order processing pipeline, for example, must
handle events (orders, payments, inventory adjustments, shipments) that do not
commute and that individually violate business rules (no shipment without
payment, no negative inventory).

We formalize the missing piece as a \emph{registry}: any authoritative
definition of validity external to the stream processor.  Registries already
exist in practice under various names---policy engines, constraint services,
schema validators, compliance controllers---but their role in convergence has
not been characterized formally.  Our model abstracts these as a tuple of
invariant predicates and a compensation operator.

This paper identifies a third regime: \textbf{normalization confluence}.
Operations may be non-commutative.  They may violate invariants.  But a
compensation operator normalizes states to validity after each operation, and
the key structural condition is that these \emph{compensated results}
commute---even though the operations themselves do not.  Under this condition,
together with well-founded compensation, Newman's Lemma guarantees that all
causality-respecting processing orders reach the same valid state.

The three regimes compare as follows:
\begin{enumerate}[leftmargin=*,nosep]
\item Operation commutativity (every state valid) $\implies$ compensation
  commutativity $\implies$ convergence.
\item Invariant confluence $\implies$ no compensation needed; in the
  sequential-application model of this paper, convergence then still
  requires the operations to commute on the reachable states.
\item Well-founded compensation and compensation commutativity $\implies$
  convergence (this paper).
\end{enumerate}
Normalization confluence occupies the design space between CRDT commutativity
(which requires operations to commute) and I-confluence (which requires
operations to preserve invariants).  It permits operations that satisfy
neither, provided their compensated results agree.

\subsection{Convergence via Normalization}

Convergence is conventionally treated as a property of \emph{operations}:
operations must be designed to commute, or designed to preserve invariants.
Our result reframes convergence as a property of the \emph{normalization
rewrite system}: the system that interleaves event application with invariant
repair.  The operations themselves need no special algebraic structure.  What
matters is that the rewrite system---encompassing both application and
compensation---is terminating and confluent.  This shifts the design burden
from operation algebra to compensation design.

\subsection{Contributions}

\begin{enumerate}[leftmargin=*,nosep]
\item \textbf{Normalization confluence regime} (\Cref{sec:model}): A
  governance rewrite system on configurations $(\sigma, B)$ pairing a state
  with a pending-event buffer, with rules for both event application and
  compensation.  Nondeterminism arises from the choice of which enabled event
  to apply next.
\item \textbf{Convergence theorem} (\Cref{sec:convergence}): Proof that
  processors consuming the same events converge to the same valid state,
  regardless of application order, under well-founded compensation (WFC) and
  compensation commutativity (CC).
\item \textbf{Tightness and exactness} (\Cref{sec:necessity}): without UBC,
  compensation depth is unbounded; a failure of CC can make different orders
  reach different valid states; and CC is the exact condition (an
  if-and-only-if) when every buffered event may fire and compensation repairs
  to validity, with a joinability form of CC exact under any enabledness.
\item \textbf{Complexity analysis} (\Cref{sec:complexity}): Proof that
  registry governance adds $O(\log n)$ overhead per event under UBC,
  preserving practical scalability.
\item \textbf{Three-regime characterization} (\Cref{sec:discussion}):
  Together with CRDTs and invariant confluence, a partition of the
  coordination-free design space into operation commutativity, invariant
  confluence, and normalization confluence.
\item \textbf{Verification calculus} (\Cref{sec:calculus}): invariant
  footprints and decomposable repair reduce the global CC obligation to
  per-event checks (strong absorption, footprint absorption) and per-pair
  commutation up to normalization, with each hypothesis shown to be needed.
  Product composition theorems enable modular verification of multi-registry
  systems.
\item \textbf{Machine-checked proof} (\Cref{sec:mechanization}): A Coq/Rocq
  mechanization of the Convergence Theorem (Newman's Lemma and the three
  critical-pair cases), of the stream-level theorem, of the exact converses,
  of the calculus, and of every worked example and counterexample, all
  axiom-free (\texttt{Print Assumptions} reports ``Closed under the global
  context''), with concrete registries witnessing that the hypotheses are
  jointly satisfiable, so no theorem is vacuous.
\end{enumerate}

%========================================
\section{Related Work}
\label{sec:related}

\paragraph{Conflict-Free Replicated Data Types.}
CRDTs~\cite{shapiro2011crdt} guarantee convergence through algebraic
structure: replicas that apply the same set of commutative, associative,
idempotent operations converge to the same state.  Our setting is strictly more
general in one dimension: we do not require operations to commute, only that
their \emph{compensated results} commute.  This permits operations that
individually violate invariants, which CRDTs cannot accommodate.  Conversely,
CRDTs require no finiteness or boundedness assumptions.  Neither approach
dominates: CRDTs are preferable when operations can be designed to commute,
while registry governance is preferable when business rules constrain valid
states in ways that preclude commutativity.

\paragraph{CALM and Coordination Avoidance.}
The CALM theorem~\cite{hellerstein2010declarative,ameloot2013relational}
establishes that monotone programs can achieve consistency without coordination.
Bailis et al.~\cite{bailis2014coordination} extend this to identify
\emph{invariant confluence}: a set of operations is invariant-confluent if all
possible orderings preserve application invariants.  When invariant confluence
holds, no compensation is needed.  Our framework addresses the complement:
operations that are \emph{not} invariant-confluent.  We prove that
compensation restores convergence under well-founded compensation and
compensation commutativity.  \Cref{sec:design-space} compares the two regimes
precisely.

\paragraph{Invariant-Repair Systems.}
Indigo~\cite{balegas2015indigo} and IPA~\cite{balegas2019ipa} also restore
application invariants over weakly consistent replicas.  They do so by a
different route: static analysis of invariant-violating operation pairs, to
which they attach application-specific conflict-resolution rules or reservation
(escrow) protocols so that every reachable state stays invariant-preserving, in
the invariant-confluence tradition.  We do not preselect operations to be
invariant-preserving.  We admit operations that genuinely violate invariants and
prove, via Newman's Lemma under compensation commutativity, that repair
reconciles all orderings to the same valid normal form.  The distinction is one
of method: reservation and invariant confluence versus confluence of a
compensating rewrite system.

\paragraph{Abstract Rewrite Systems.}
Newman's Lemma~\cite{newman1942} states that a terminating abstract rewrite
system is confluent if and only if it is locally confluent.  The lemma has been
applied extensively in term rewriting~\cite{baader1998term} and lambda
calculus.  Our contribution is its application to a rewrite system with
\emph{genuine nondeterminism}: the choice of which event to apply next from a
buffer of enabled events.  Local confluence is nontrivial because different
application orders produce different intermediate states; our compensation
commutativity axiom is the condition that closes the resulting critical pairs.

\paragraph{Rewriting Logic and Coherence.}
Meseguer's rewriting logic~\cite{meseguer1992} models concurrent systems as
rewrite theories, and coherence checking (the Church-Rosser and coherence
properties of rewrite theories modulo equations~\cite{duran2012coherence},
supported by tools such as Maude) is the closest methodological ancestor of our
build-time convergence check.  We instantiate this lineage for one specific
rewrite system: events interleaved with invariant compensation, where
compensation commutativity plays the role of the coherence condition that closes
critical pairs.

\paragraph{Self-Stabilization.}
Self-stabilizing systems~\cite{dijkstra1974} converge to a legitimate state from
an arbitrary starting configuration.  Normalization confluence shares the
converge-to-valid premise but strengthens the conclusion to agreement: not
merely that each replica reaches some valid state, but that all replicas reach
the \emph{same} valid state under reordering.

\paragraph{Eventual Consistency.}
Classical eventual consistency~\cite{vogels2009eventually} guarantees that
replicas converge given sufficient time without new updates.  It does not
address whether converged states satisfy application-level invariants---only
that replicas agree on the same value.  Our work guarantees convergence to
\emph{valid} states under continuous event arrival.

\paragraph{Stream Processing Theory.}
Formal models of stream processing include synchronous
dataflow~\cite{lee1987synchronous}, Kahn process
networks~\cite{kahn1974semantics}, and the windowed computation models
underlying modern engines~\cite{akidau2015dataflow}.  These focus on
operational semantics---what a stream processor computes---rather than semantic
validity of computed states.  Our registry layer sits above any of these
operational models.

\paragraph{Saga Pattern and Compensating Transactions.}
The saga pattern~\cite{garcia1987sagas} decomposes long-lived transactions into
sequences of local transactions with compensating actions for rollback.  Our
compensation operator generalizes this: rather than undoing individual
operations, it repairs arbitrary invariant violations in the current state.
The convergence guarantee we prove---that different interleaving orders of
application and compensation reach the same valid state---has no analogue in
the saga literature.

\paragraph{Mechanized Convergence Proofs.}
Gomes et al.~\cite{gomes2017verifying} give a machine-checked proof of strong
eventual consistency for CRDTs in Isabelle/HOL, network model included.  Their
result covers the commutative case; the mechanization accompanying this work
(\Cref{sec:mechanization}) covers convergence under non-commutative,
invariant-violating operations reconciled by compensation.  Both rest on an
assumption-free, machine-checked convergence theorem rather than a paper proof.

%========================================
\section{Registry-Governed Stream Model}
\label{sec:model}

\subsection{Registries and Semantic Projection}

\begin{definition}[Registry]
\label{def:registry}
A \emph{registry} is a tuple $R = (\Sigma_R, I_R, V_R, \rho_R)$ where:
\begin{itemize}[leftmargin=*,nosep]
\item $\Sigma_R$ is a set of \emph{semantic states},
\item $I_R = \{\psi_1, \ldots, \psi_k\}$ is a finite set of \emph{invariant
  predicates} $\psi_i: \Sigma_R \to \{\top, \bot\}$,
\item $V_R: \Sigma_R \to \{\top, \bot\}$ is the \emph{validation function}
  defined by $V_R(\sigma) = \bigwedge_{i=1}^k \psi_i(\sigma)$,
\item $\rho_R: \Sigma_R \to \Sigma_R$ is a \emph{compensation operator}
  satisfying: if $V_R(\sigma) = \bot$, then $\Phi(\rho_R(\sigma)) < \Phi(\sigma)$
  for a well-founded measure $\Phi$ (defined below); if
  $V_R(\sigma) = \top$, then $\rho_R(\sigma) = \sigma$.
\end{itemize}
\end{definition}

\begin{definition}[Semantic Projection]
\label{def:projection}
A registry is defined over a semantic state space $\Sigma_R$ together with
a \emph{projection} $\pi: \Sigma_{\mathrm{raw}} \to \Sigma_R$ from
unbounded raw states to semantic states.  All invariants, validation, and
compensation are defined on $\Sigma_R$.  Two raw states $s_1, s_2$ are
\emph{registry-indistinguishable} if $\pi(s_1) = \pi(s_2)$.
The application function $\mathrm{apply}: \mathcal{E} \times \Sigma_R \to \Sigma_R$
is understood as ``apply then project'': given a raw-level application
$\mathrm{apply}_{\mathrm{raw}}: \mathcal{E} \times \Sigma_{\mathrm{raw}} \to \Sigma_{\mathrm{raw}}$,
the semantic application is
$\mathrm{apply}(e, \sigma) = \pi(\mathrm{apply}_{\mathrm{raw}}(e, s))$ for
any $s$ with $\pi(s) = \sigma$.  Well-definedness requires that the projected
result is independent of the representative~$s$.
\end{definition}

\begin{remark}[Scope of Finiteness]
\label{rem:scope}
Finiteness is required of the \emph{semantic} state space $\Sigma_R$, not the
raw data.  An order-processing registry tracking status across states
\emph{pending}, \emph{paid}, \emph{shipped}, \emph{delivered},
\emph{cancelled} has $|\Sigma_R| = 5$ regardless of how many bytes the
underlying order record occupies.  More generally, any domain whose logic is
captured by $k$ boolean predicates has $|\Sigma_R| \leq 2^k$: the semantic
states are the equivalence classes under the invariant vector
$(\psi_1(\sigma), \ldots, \psi_k(\sigma))$.  Business domains are
overwhelmingly of this form.
\end{remark}

\begin{definition}[Registry Equivalence]
\label{def:equiv}
Two states $\sigma_1, \sigma_2 \in \Sigma_R$ are \emph{registry equivalent},
written $\sigma_1 \approx_R \sigma_2$, if they satisfy the same invariants:
$\forall \psi \in I_R: \psi(\sigma_1) = \psi(\sigma_2)$.
\end{definition}

\begin{lemma}
\label{lem:equiv-relation}
Registry equivalence $\approx_R$ is an equivalence relation on~$\Sigma_R$.
\end{lemma}
\begin{proof}
Each condition $\psi(\sigma_1) = \psi(\sigma_2)$ is reflexive, symmetric, and
transitive.  A conjunction of equivalence relations is an equivalence relation.
\end{proof}

\subsection{Well-Founded and Uniformly Bounded Compensation}

Termination and complexity require two properties of the compensation
operator, which we state separately because they play different roles.

\begin{axiom}[Well-Founded Compensation (WFC)]
\label{ax:wfc}
There exists a \emph{compensation measure}
$\Phi: \Sigma_R \to \mathbb{N}$ such that every compensation step strictly
decreases $\Phi$: if $V_R(\sigma) = \bot$, then
$\Phi(\rho_R(\sigma)) < \Phi(\sigma)$.
\end{axiom}

WFC guarantees that compensation terminates from any state: iterated
application of~$\rho_R$ reaches a valid state in finitely many steps.

\begin{axiom}[Uniformly Bounded Compensation (UBC)]
\label{ax:ubc}
The compensation measure $\Phi$ from \Cref{ax:wfc} is uniformly bounded:
there exists $M \in \mathbb{N}$ such that $\Phi(\sigma) \leq M$ for all
$\sigma \in \Sigma_R$.
\end{axiom}

UBC strengthens WFC by bounding compensation depth \emph{uniformly} across
all states.  WFC alone suffices for termination and for defining iterated
compensation~$\rho_R^*$.  UBC is additionally needed for the uniform
step bound in the termination lemma and the constant per-event compensation
cost in the complexity analysis.

\begin{lemma}[Finite State Spaces Imply UBC]
\label{lem:finite-implies-ubc}
If $|\Sigma_R| < \infty$ and WFC holds, then UBC holds with
$M = \max_{\sigma \in \Sigma_R} \Phi(\sigma)$.
\end{lemma}
\begin{proof}
The maximum exists because $\Sigma_R$ is finite.
\end{proof}

\begin{remark}
WFC is the load-bearing condition for termination; UBC is the load-bearing
condition for uniform step bounds and per-event complexity.  Finiteness of
$\Sigma_R$ is a convenient sufficient condition for both, but UBC does not
require full finiteness of $\Sigma_R$; it requires only bounded compensation
height.  The counterexample
in \Cref{sec:necessity} satisfies WFC (compensation terminates for every fixed
event set) but violates UBC (compensation depth grows without bound as event
sets grow).
\end{remark}

\subsection{Events and Streams}

\begin{definition}[Event Stream]
\label{def:stream}
An \emph{event stream} $\mathcal{S} = \langle e_1, e_2, \ldots \rangle$ is a
(possibly infinite) sequence of events drawn from event space $\mathcal{E}$.
Each event $e_i$ carries a unique identifier $\mathrm{id}(e_i)$, a timestamp
$\tau(e_i)$, and a set of causal dependencies
$\mathrm{deps}(e_i) \subseteq \{e_1, \ldots, e_{i-1}\}$.  An
\emph{application function}
$\mathrm{apply}: \mathcal{E} \times \Sigma_R \to \Sigma_R$ maps each event
and current state to a new state.
\end{definition}

\begin{definition}[Causality-Respecting Order]
\label{def:causal-order}
An ordering $\pi$ of a finite event set $E$ \emph{respects causality} if
whenever $e_i \in \mathrm{deps}(e_j)$, event $e_i$ precedes $e_j$
under~$\pi$.  Multiple causality-respecting orders may exist for the same
event set.
\end{definition}

\subsection{Governance Rewrite System}

We now define the rewrite system that is the core technical device of this
paper.  Unlike the approach of fixing a canonical event order---which would make
agreement trivial by determinism---our rewrite system admits genuine
nondeterminism in the choice of which event to apply next.

\begin{assumption}[Causal Closure]
\label{as:closure}
The received event set $E$ is causally closed: for every $e \in E$,
$\mathrm{deps}(e) \subseteq E$.
\end{assumption}

Under causal closure, every dependency of every event in~$E$ is either still
pending in~$B$ or has already been applied (removed from~$B$).  This makes the
enabled predicate below well-defined.

\begin{definition}[Configuration]
\label{def:config}
A \emph{configuration} is a pair $(\sigma, B)$ where $\sigma \in \Sigma_R$ is
a semantic state and $B \subseteq E$ is a set of events that have been
received but not yet applied.  An event $e \in B$ is \emph{enabled} if
$\mathrm{deps}(e) \cap B = \varnothing$ (all causal dependencies of $e$ have
already been applied).
\end{definition}

\begin{definition}[Governance Rewrite System]
\label{def:grs}
Fix a finite set of received events~$E$ and initial state~$\sigma_0$.  The
\emph{governance rewrite system} $\mathcal{G} = (\mathcal{C}, \to)$ operates
on configurations $(\sigma, B)$ with $\sigma \in \Sigma_R$ and
$B \subseteq E$.  The rewrite relation $\to$ is the union of:
\begin{enumerate}[leftmargin=*,nosep]
\item \textbf{Apply step:}
  $(\sigma, B) \to (\mathrm{apply}(e, \sigma), B \setminus \{e\})$ whenever
  $e \in B$ is enabled.
\item \textbf{Compensation step:}
  $(\sigma, B) \to (\rho_R(\sigma), B)$ whenever $V_R(\sigma) = \bot$.
\end{enumerate}
A configuration is in \emph{normal form} if $B = \varnothing$ and
$V_R(\sigma) = \top$.
\end{definition}

The nondeterminism is genuine: when multiple events are enabled and the state
is invalid, a configuration may step via different apply rules (choosing
different events) or via the compensation rule.  Confluence of this system
means that the choice does not affect the final result.

\subsection{Compensation Commutativity}

Local confluence of $\mathcal{G}$ is nontrivial.  Two critical pair types
arise: (1)~two apply steps choosing different events, and (2)~an apply step
versus a compensation step.  The following axiom ensures both types close.

\begin{definition}[Iterated Compensation]
\label{def:rhostar}
Under WFC (\Cref{ax:wfc}), define $\rho_R^*(\sigma)$ to be the unique state
obtained by iterating $\rho_R$ from~$\sigma$ until reaching validity:
$\rho_R^*(\sigma) = \rho_R^{m}(\sigma)$ for the least $m$ such that
$V_R(\rho_R^{m}(\sigma)) = \top$.  Existence and finiteness of~$m$ follow from
WFC (indeed $m \le \Phi(\sigma)$); uniqueness follows from determinism
of~$\rho_R$.  Because $V_R$ is Boolean-valued, $\rho_R^*$ is computable, and it
does not depend on which measure witnesses WFC.
\end{definition}

\begin{axiom}[Compensation Commutativity (CC)]
\label{ax:cc}
The registry $R$ satisfies:
\begin{enumerate}[leftmargin=*,nosep]
\item[\textbf{(CC1)}] \textbf{Order independence.}\; For any
  $\sigma \in \Sigma_R$ and events $e_1, e_2$ that are causally independent
  and simultaneously enabled in some configuration $(\sigma, B)$:
  \[
  \rho_R^*(\mathrm{apply}(e_2,\, \rho_R^*(\mathrm{apply}(e_1, \sigma))))
  \;=\;
  \rho_R^*(\mathrm{apply}(e_1,\, \rho_R^*(\mathrm{apply}(e_2, \sigma)))).
  \]
\item[\textbf{(CC2)}] \textbf{Compensation absorption.}\; For any
  $\sigma \in \Sigma_R$ with $V_R(\sigma) = \bot$ and any enabled event~$e$:
  \[
  \rho_R^*(\mathrm{apply}(e, \sigma))
  \;=\;
  \rho_R^*(\mathrm{apply}(e, \rho_R(\sigma))).
  \]
\end{enumerate}
\end{axiom}

\begin{remark}[CRDT Commutativity]
\label{rem:crdt-comparison}
CRDT commutativity requires
$\mathrm{apply}(e_1, \mathrm{apply}(e_2, \sigma)) =
\mathrm{apply}(e_2, \mathrm{apply}(e_1, \sigma))$ for all~$\sigma$.  This
implies CC1 (take $\rho_R^* = \mathrm{id}$ when all states are valid) and
trivially satisfies CC2 (when all states are valid, compensation never fires).
Compensation commutativity is strictly weaker: it permits operations that
individually violate invariants, requiring only that the \emph{repaired}
results commute.  The gap between operation commutativity and compensation
commutativity is precisely the space our result occupies.
\end{remark}

\begin{remark}[Compensation Absorption]
\label{rem:absorption}
Condition CC2 states that compensating an invalid state before applying an
event does not change the final compensated result.  Intuitively, the
repair that compensation performs ``would have been performed anyway'' during
the post-application compensation phase.  CC2 holds whenever compensation
repairs invariant violations in a way that is determined by the violation
pattern rather than the path taken to reach the violation.  By induction on
the number of compensation steps, CC2 generalizes:
$\rho_R^*(\mathrm{apply}(e, \sigma)) =
\rho_R^*(\mathrm{apply}(e, \rho_R^*(\sigma)))$.
In rewrite-theoretic terms, CC2 ensures that the normalizer is stable under
interleaving with event application: normalizing ``early'' (before an event)
versus ``late'' (after) does not change the final normal form.
\end{remark}

\subsection{Stream Processors}

\begin{definition}[Stream Processor]
\label{def:processor}
A \emph{stream processor} $P$ for stream $\mathcal{S}$ under registry $R$ is a
process that, at each time $t$, has received some subset of events
$E_P(t) \subseteq \mathcal{S}$ and maintains state $\Psi_P(t) \in \Sigma_R$.
We require:
\begin{enumerate}[leftmargin=*,nosep]
\item \textbf{Eventual delivery}: For every event $e_i \in \mathcal{S}$, there
  exists a time $t_i$ such that $e_i \in E_P(t)$ for all $t \geq t_i$.
\item \textbf{Causality-respecting application}: The processor applies events
  in some order that respects causality (\Cref{def:causal-order}), buffering
  events whose causal dependencies have not yet been received.  The processor
  is free to choose any causality-respecting order.
\item \textbf{Incremental compensation}: After applying each event, the
  processor applies $\rho_R$ until validity is restored.
\end{enumerate}
The processor's computation on event set $E_P(t)$ is a reduction sequence in
the governance rewrite system $\mathcal{G}$ starting from
$(\sigma_0, E_P(t))$.  The processor has \emph{finished processing} $E_P(t)$
at time~$t$ if its configuration at~$t$ is a normal form of that reduction.
It is \emph{fair} if, while no new events arrive, it keeps reducing, and makes
progress whenever it has not finished.
\end{definition}

A received set that is not yet causally closed (an event has arrived before
one of its dependencies) is never finished in this sense: the waiting event
stays in the buffer, so the configuration is not a normal form.  Different processors may apply the same events in different
causality-respecting orders.  The convergence theorem guarantees that, once
they have finished processing the same events, they hold the same valid state.

%========================================
\section{Worked Example: Order Fulfillment}
\label{sec:example}

We instantiate the framework on an order fulfillment pipeline, verify the
structural conditions, and trace two processor orderings to the same valid
state.

\subsection{Registry Definition}

An order has a \emph{stage} in $\{\textit{pending}, \textit{approved},
\textit{held}\}$ and a \emph{balance} in $\{0, 1, 2, 3\}$.  The semantic
state space is $\Sigma_R = \{\textit{pending}, \textit{approved},
\textit{held}\} \times \{0, 1, 2, 3\}$, with $|\Sigma_R| = 12$ and initial
state $\sigma_0 = (\textit{pending}, 0)$.

A single invariant governs fulfillment:
\[
\psi(s, b) \;=\; \top \;\iff\; (s = \textit{approved} \implies b \geq 2).
\]
Approval requires a balance of at least~2.  The only invalid states are
$(\textit{approved}, 0)$ and $(\textit{approved}, 1)$.

\paragraph{Compensation.}
When an order is approved with insufficient balance, it is placed on hold:
\[
\rho(\textit{approved}, b) = (\textit{held}, b) \quad \text{for } b < 2.
\]
All other states are valid, and $\rho$ acts as the identity.  The
\textit{held} stage preserves the intent to approve: the system records that
approval was requested but cannot yet proceed.  This is a deliberate design
choice---a naive alternative that simply cancels the approval (reverting to
\textit{pending}) would violate CC, as we discuss below.

\paragraph{Measures.}
$\Phi(s, b) = 1$ if $s = \textit{approved}$ and $b < 2$, otherwise $0$.
WFC holds ($\rho$ strictly decreases $\Phi$ on invalid states), and UBC holds
with $M = 1$.

\subsection{Events}

Two causally independent events:
\begin{itemize}[leftmargin=*,nosep]
\item $e_{\mathrm{credit}}$: Adds 1 to the balance (capped at~3).  If the
  order is on hold and the new balance reaches~2, the order is automatically
  approved:
  \[
  \mathrm{apply}(e_{\mathrm{credit}}, (s, b)) =
  \begin{cases}
  (\textit{approved}, \min(b{+}1, 3)) & \text{if } s = \textit{held} \text{ and } \min(b{+}1,3) \geq 2, \\
  (s, \min(b{+}1, 3)) & \text{otherwise.}
  \end{cases}
  \]
\item $e_{\mathrm{approve}}$: Sets the stage to \textit{approved}:
  $\mathrm{apply}(e_{\mathrm{approve}}, (s, b)) = (\textit{approved}, b)$.
\end{itemize}

\subsection{Verifying CC}

\paragraph{CC1 (Order Independence).}
We verify CC1 for $(e_{\mathrm{approve}}, e_{\mathrm{credit}})$ at two
representative starting states.

\emph{From $(\textit{pending}, 0)$}:
\begin{align*}
&\text{Path 1 ($e_{\mathrm{approve}}$ first):} \\
&\quad (\textit{pending}, 0) \xrightarrow{e_{\mathrm{approve}}}
  (\textit{approved}, 0) \xrightarrow{\rho^*} (\textit{held}, 0)
  \xrightarrow{e_{\mathrm{credit}}} (\textit{held}, 1).
  \quad\text{Result: } (\textit{held}, 1). \\[4pt]
&\text{Path 2 ($e_{\mathrm{credit}}$ first):} \\
&\quad (\textit{pending}, 0) \xrightarrow{e_{\mathrm{credit}}}
  (\textit{pending}, 1) \xrightarrow{e_{\mathrm{approve}}}
  (\textit{approved}, 1) \xrightarrow{\rho^*} (\textit{held}, 1).
  \quad\text{Result: } (\textit{held}, 1). \;\checkmark
\end{align*}

In Path~1, approval fires first, violates the invariant (balance~0 $< 2$),
and compensation places the order on hold.  The subsequent credit increments
the balance but does not yet reach the threshold.  In Path~2, credit arrives
first (no violation), then approval violates the invariant and compensation
again produces the held state.  Both paths reach $(\textit{held}, 1)$.

\emph{From $(\textit{pending}, 1)$}:
\begin{align*}
&\text{Path 1:}\;
  (\textit{pending}, 1) \xrightarrow{e_{\mathrm{approve}}}
  (\textit{approved}, 1) \xrightarrow{\rho^*} (\textit{held}, 1)
  \xrightarrow{e_{\mathrm{credit}}} (\textit{approved}, 2).
  \quad\text{Result: } (\textit{approved}, 2). \\
&\text{Path 2:}\;
  (\textit{pending}, 1) \xrightarrow{e_{\mathrm{credit}}}
  (\textit{pending}, 2) \xrightarrow{e_{\mathrm{approve}}}
  (\textit{approved}, 2).
  \quad\text{Result: } (\textit{approved}, 2). \;\checkmark
\end{align*}

In Path~1, approval again violates the invariant; compensation holds; the
subsequent credit reaches balance~2, triggering auto-approval.  In Path~2, the
credit arrives first, and approval at balance~2 is valid.  Both reach
$(\textit{approved}, 2)$: the order is approved with sufficient funds.

Since $|\Sigma_R| = 12$, CC1 can be verified exhaustively; it holds for every
pair of events at all 12 states, invalid ones included.  All cases follow the
same pattern: compensation preserves the approval intent in the held state,
and subsequent credits eventually satisfy the balance threshold.

\paragraph{CC2 (Compensation Absorption).}
From the invalid state $(\textit{approved}, 0)$ with event
$e_{\mathrm{credit}}$:
\begin{align*}
\text{LHS:}\;\; & \rho^*(\mathrm{apply}(e_{\mathrm{credit}},\,
  (\textit{approved}, 0)))
  = \rho^*((\textit{approved}, 1))
  = (\textit{held}, 1). \\
\text{RHS:}\;\; & \rho^*(\mathrm{apply}(e_{\mathrm{credit}},\,
  \rho(\textit{approved}, 0)))
  = \rho^*(\mathrm{apply}(e_{\mathrm{credit}},\,
  (\textit{held}, 0)))
  = \rho^*((\textit{held}, 1))
  = (\textit{held}, 1). \;\checkmark
\end{align*}

Compensating before the credit (putting the order on hold first) produces the
same result as applying the credit to the invalid state and then compensating.
CC2 holds because the held state and the approved-but-invalid state respond
identically to subsequent credits after compensation.

\subsection{Two Processors, Three Events}

Consider the stream $\mathcal{S} = \langle e_{\mathrm{approve}},
e_{\mathrm{credit}}, e_{\mathrm{credit}} \rangle$ (approval followed by two
credits).  Two processors $P_1, P_2$ receive the events in different orders:

\medskip\noindent
\textbf{Processor $P_1$} (approve, credit, credit):
\[
(\textit{pending}, 0)
\xrightarrow{e_{\mathrm{approve}}} (\textit{approved}, 0)
\xrightarrow{\rho^*} (\textit{held}, 0)
\xrightarrow{e_{\mathrm{credit}}} (\textit{held}, 1)
\xrightarrow{e_{\mathrm{credit}}} (\textit{approved}, 2).
\]

\noindent
\textbf{Processor $P_2$} (credit, credit, approve):
\[
(\textit{pending}, 0)
\xrightarrow{e_{\mathrm{credit}}} (\textit{pending}, 1)
\xrightarrow{e_{\mathrm{credit}}} (\textit{pending}, 2)
\xrightarrow{e_{\mathrm{approve}}} (\textit{approved}, 2).
\]

\noindent Both processors arrive at $(\textit{approved}, 2)$.  $P_1$ encountered
an invariant violation (approval without sufficient balance) and passed
through the held state; $P_2$ never violated the invariant.  Compensation
commutativity guarantees agreement regardless of the path.

\subsection{Why Compensation Design Matters}

The held stage is essential.  A naive compensation that simply reverts
approval to \textit{pending}, that is,
$\rho(\textit{approved}, b) = (\textit{pending}, b)$ for $b < 2$, satisfies WFC
and UBC but violates CC1.  The witness is the state $(\textit{pending}, 1)$:
approval then credit gives $\rho(\textit{approved}, 1) = (\textit{pending}, 1)$
and then $(\textit{pending}, 2)$ (intent to approve lost), while credit then
approval gives $(\textit{pending}, 2)$ and then $(\textit{approved}, 2)$, which
is valid.  So
\begin{align*}
\rho^*(\mathrm{apply}(e_{\mathrm{credit}},
  \rho^*(\mathrm{apply}(e_{\mathrm{approve}}, (\textit{pending}, 1)))))
&= (\textit{pending}, 2), \\
\rho^*(\mathrm{apply}(e_{\mathrm{approve}},
  \rho^*(\mathrm{apply}(e_{\mathrm{credit}}, (\textit{pending}, 1)))))
&= (\textit{approved}, 2) \neq (\textit{pending}, 2).
\end{align*}
On the stream $\langle e_{\mathrm{approve}}, e_{\mathrm{credit}},
e_{\mathrm{credit}} \rangle$ below, $P_1$ would end in $(\textit{pending}, 2)$
and $P_2$ in $(\textit{approved}, 2)$.  (At $(\textit{pending}, 0)$ the two
orders happen to agree: both give $(\textit{pending}, 1)$.)  The naive revert
also violates CC2: at the invalid state $(\textit{approved}, 1)$,
\begin{align*}
\rho^*(\mathrm{apply}(e_{\mathrm{credit}}, (\textit{approved}, 1)))
&= (\textit{approved}, 2), \\
\rho^*(\mathrm{apply}(e_{\mathrm{credit}}, \rho(\textit{approved}, 1)))
&= (\textit{pending}, 2).
\end{align*}
The held stage preserves the approval intent so that
subsequent credits can trigger auto-approval, ensuring that different orderings
converge.

This illustrates the practical content of the CC conditions: they constrain
\emph{compensation design}, not operation design.  What the held stage provides
is \emph{absorption}: compensating before a credit or after it gives the same
normalized result (CC2), and the two events commute once results are
normalized.  \Cref{sec:calculus} shows that these per-event and per-pair facts
give CC1 at every state, and that under absorption CC1 is equivalent to
commutation up to normalization (\Cref{thm:footprint-cc1}).

%========================================
\section{Convergence Theorem}
\label{sec:convergence}

\subsection{Termination}

\begin{lemma}[Termination]
\label{lem:termination}
Under WFC (\Cref{ax:wfc}), the governance rewrite system $\mathcal{G}$ is
terminating: every reduction sequence from $(\sigma_0, E)$ is finite.
Under UBC (\Cref{ax:ubc}), the total number of steps is bounded by
$|E| + (|E| + 1) \cdot M$.
\end{lemma}
\begin{proof}
Define the measure $\mu(\sigma, B) = (|B|, \Phi(\sigma))$ taking values in
$\mathbb{N} \times \mathbb{N}$ under the lexicographic order, which is
well-founded.

An apply step sends $(\sigma, B)$ to $(\mathrm{apply}(e, \sigma), B \setminus \{e\})$.
The first component decreases: $|B \setminus \{e\}| = |B| - 1$.  Therefore
$\mu$ strictly decreases regardless of the second component.

A compensation step sends $(\sigma, B)$ to $(\rho_R(\sigma), B)$.  The first
component is unchanged; the second strictly decreases by WFC.
Therefore $\mu$ strictly decreases.

Since $\mu$ maps into a well-ordered set and strictly decreases at every step,
every reduction sequence is finite.  Under UBC, $\Phi(\sigma) \leq M$ for all
$\sigma$, so at most $M$ compensation steps can occur between apply steps
(including before the first apply), giving the uniform bound
$|E| + (|E| + 1) \cdot M$.
\end{proof}

The bound is tight: a registry with $M = 1$ and a single event has a
reduction of $3 = |E| + (|E|+1) M$ steps (compensate, apply, compensate).

\subsection{Local Confluence}

\begin{lemma}[Local Confluence]
\label{lem:local-confluence}
Under CC (\Cref{ax:cc}), the governance rewrite system $\mathcal{G}$ is
locally confluent: if $(\sigma, B) \to X_1$ and $(\sigma, B) \to X_2$,
then there exists a configuration~$Z$ with $X_1 \to^* Z$ and
$X_2 \to^* Z$.
\end{lemma}
\begin{proof}
Three cases arise from the two rewrite rules.

\medskip\noindent\textbf{Case 1: Two apply steps.}\;
$(\sigma, B) \to (\mathrm{apply}(e_1, \sigma), B \setminus \{e_1\})
= X_1$ and
$(\sigma, B) \to (\mathrm{apply}(e_2, \sigma), B \setminus \{e_2\})
= X_2$, where $e_1 \neq e_2$ are both enabled.

Since $e_1$ and $e_2$ are both enabled, their dependencies are disjoint
from~$B$.  By \Cref{as:closure}, $e_1, e_2 \in E$, and since $e_1 \in B$ we
have $e_1 \notin \mathrm{deps}(e_2)$ (otherwise
$\mathrm{deps}(e_2) \cap B \neq \varnothing$, contradicting enabledness of
$e_2$), and symmetrically $e_2 \notin \mathrm{deps}(e_1)$.  Therefore $e_1$
and $e_2$ are causally independent.

From $X_1$: compensate to validity, obtaining
$(\rho_R^*(\mathrm{apply}(e_1, \sigma)),\, B \setminus \{e_1\})$; then
apply~$e_2$ (still enabled since its dependencies remain outside
$B \setminus \{e_1\}$) and compensate, reaching
\[
Z_1 = \bigl(\rho_R^*(\mathrm{apply}(e_2,\,
  \rho_R^*(\mathrm{apply}(e_1, \sigma)))),\;
  B \setminus \{e_1, e_2\}\bigr).
\]
From $X_2$: symmetrically, reach
\[
Z_2 = \bigl(\rho_R^*(\mathrm{apply}(e_1,\,
  \rho_R^*(\mathrm{apply}(e_2, \sigma)))),\;
  B \setminus \{e_1, e_2\}\bigr).
\]
By CC1, the state components of $Z_1$ and $Z_2$ are equal.  Take $Z = Z_1 = Z_2$.

\medskip\noindent\textbf{Case 2: Apply vs.\ compensation.}\;
$V_R(\sigma) = \bot$ and $e$ is enabled in~$B$.
$(\sigma, B) \to (\mathrm{apply}(e, \sigma), B \setminus \{e\}) = X_1$
and
$(\sigma, B) \to (\rho_R(\sigma), B) = X_2$.

From $X_1$: compensate to validity, reaching
$(\rho_R^*(\mathrm{apply}(e, \sigma)),\, B \setminus \{e\})$.

From $X_2 = (\rho_R(\sigma), B)$: apply~$e$ (still enabled, since
compensation does not change~$B$) to get
$(\mathrm{apply}(e, \rho_R(\sigma)),\, B \setminus \{e\})$; then
compensate to validity, reaching
$(\rho_R^*(\mathrm{apply}(e, \rho_R(\sigma))),\, B \setminus \{e\})$.

By CC2, $\rho_R^*(\mathrm{apply}(e, \sigma)) =
\rho_R^*(\mathrm{apply}(e, \rho_R(\sigma)))$.  The buffer components
are equal.  Take $Z = (\rho_R^*(\mathrm{apply}(e, \sigma)),\, B \setminus
\{e\})$.

\medskip\noindent\textbf{Case 3: Two compensation steps.}\;
Both rewrites apply $\rho_R$ to the same state, producing the same
successor $(\rho_R(\sigma), B)$.  Take $Z = (\rho_R(\sigma), B)$.
\end{proof}

\subsection{Main Result}

\begin{corollary}[Unique Normal Forms]
\label{cor:unique-nf}
Under WFC and CC, every configuration $(\sigma_0, E)$ has a unique normal
form.
\end{corollary}
\begin{proof}
By \Cref{lem:termination}, $\mathcal{G}$ is terminating.  By
\Cref{lem:local-confluence}, $\mathcal{G}$ is locally confluent.  By Newman's
Lemma~\cite{newman1942}, $\mathcal{G}$ is globally confluent.  A terminating,
confluent rewrite system has unique normal forms.
\end{proof}

\begin{remark}[Machine-checked]
\label{sec:mechanization}
This result and the rest of the paper's theory are mechanized in Coq/Rocq
(checked with Rocq~9.3; CI also runs Coq~8.18 and~8.20).  Newman's Lemma is
proved from scratch (no external libraries) as \emph{strongly normalizing
$+$ locally confluent $\implies$ confluent}, and the Governance Rewrite System
is formalized with its two rules; the lexicographic measure
$(|B|, \Phi(\sigma))$ discharges termination and the three critical-pair cases
(closed by CC1, CC2, and determinism of~$\rho_R$) discharge local confluence,
yielding confluence and unique normal forms.  The gate script
\texttt{coq/verify.sh} checks 715 named results, and \texttt{Print
Assumptions} reports ``Closed under the global context'' (no axioms, no
admitted lemmas) for every one.  The results this paper relies on, by
section:
\begin{itemize}[leftmargin=*,nosep]
\item \emph{Model and convergence.}  $\rho_R^*$ is constructed from WFC, as
  in \Cref{def:rhostar} (\coq{base_def_rhostar},
  \coq{base_def_rhostar_least_unique}, \coq{rho_star_canonical}; properties
  \coq{rho_star_valid}, \coq{rho_star_fix}, \coq{rho_star_idem}), so
  \Cref{cor:unique-nf} holds with WFC, CC1 on distinct, independent,
  co-enabled pairs, and CC2 as its only hypotheses, with the
  dependency-based enabledness of \Cref{def:config}
  (\coq{paper_causal_governance_unique_normal_forms},
  \coq{paper_governance_unique_normal_forms}; also
  \coq{governance_unique_normal_forms},
  \coq{causal_governance_unique_normal_forms}).  Termination and its bound:
  \coq{base_lem_termination_bound}, attained by
  \coq{lv_termination_bound_tight}; \Cref{lem:finite-implies-ubc}:
  \coq{base_lem_finite_implies_ubc}; \Cref{rem:absorption}:
  \coq{base_rem_absorption}.
\item \emph{Stream level} (\Cref{def:processor}, \Cref{thm:convergence}):
  \coq{base_thm_convergence}, \coq{stream_convergence},
  \coq{stream_agreement}, \coq{base_rem_set_function},
  \coq{base_cor_quiescent}; counterexamples to the version-1 statements
  \coq{base_thm_convergence_c_counterexample},
  \coq{base_thm_convergence_transient_counterexample}.
\item \emph{Domain independence} (\Cref{cor:infinite}): WFC over any
  well-founded order (\coq{governance_wf_confluent},
  \coq{causal_governance_wf_confluent}, \coq{base_cor_infinite}), with an
  instance on~$\mathbb{Z}$ (\coq{zw_confluent}, \coq{zw_stream_agree}).
\item \emph{Tightness and exactness} (\Cref{sec:necessity}):
  \coq{thm_necessity}, \coq{ri_depth_exact}, \coq{ri_no_ubc},
  \coq{ri_what_fails}, \coq{ri_cc_any_extension}; \coq{prop_cc_necessary};
  the exact converses \coq{cc_exact_from} (with $\rho_R^*$ constructed:
  \coq{wfc_cc_exact_from}) and \coq{jc_exact}, and the qualifiers
  \coq{masked_cc1}, \coq{rho_star_qualifier}.
\item \emph{Worked example} (\Cref{sec:example}): \coq{of_registry},
  \coq{of_cc1}, \coq{of_cc2}, \coq{of_paper_traces}, \coq{of_processors};
  the naive revert: \coq{naive_cc1_fails}, \coq{naive_cc2_fails},
  \coq{naive_stream_diverges}, \coq{naive_paper_witness_fails}.
\item \emph{Complexity} (\Cref{thm:complexity}, model-level step bound):
  \coq{base_thm_complexity_per_event}, \coq{base_thm_complexity_comp_total}.
\item \emph{Calculus} (\Cref{sec:calculus}): \coq{strong_absorption_iff_cc2},
  \coq{base_lem_repair_commute}, \coq{base_lem_repair_idempotent}, the
  corrected footprint theorem and its parts (\coq{calc_footprint_cc1_iff},
  \coq{calc_footprint_cc1_valid_iff}, \coq{calc_footprint_cc1},
  \coq{calc_fp_absorb_iff_sa}, \coq{calc_cc1_iff_commute_under_cc2}), the
  counterexamples to the version-1 statement
  (\coq{base_thm_footprint_cc1_refuted},
  \coq{base_thm_footprint_cc1_raw_refuted}) and to dropping each hypothesis
  (\coq{calc_valid_needs_commute}, \coq{calc_valid_needs_disjoint},
  \coq{calc_valid_needs_repair_local}, \coq{calc_all_needs_commute},
  \coq{calc_all_needs_absorb}, \coq{calc_all_needs_repair_local}), the
  canonical-state pattern (\coq{calc_canonical_pattern},
  \coq{base_pattern1_cc1_refuted}), and product lifting
  (\coq{base_thm_product}, \coq{base_thm_product_decomposable}).
\end{itemize}
Each abstract theorem comes with a concrete registry discharging all of its
hypotheses with genuine nondeterminism (for example \coq{example_confluent},
\coq{calc_clamp_instance}, \coq{calc_order_cc1}), so no theorem is vacuous.
The soundness of the reference implementation's own certification strategy
is mechanized as well: events whose state-variable footprints (read and
write sets) are disjoint commute (\coq{disjoint_events_commute}; this concerns
state variables and is a different statement from the invariant footprints of
\Cref{sec:calculus}), and a repair strictly decreasing a potential is
strongly normalizing (\coq{repair_terminates}).  The development is
furthermore paired with two runnable checkers extracted from it, each proven
axiom-free: one certifies that a compiled machine's emitted step tables
commute and stay in range (\coq{check_tables_converges}), and one recomputes
convergence directly from the machine's rules by evaluating their expression
trees, trusting neither the implementation's enumeration nor its
normalization (\coq{checkBuild_converges}).  Because the checking logic is
extracted from the proof rather than reimplemented by hand, the two serve as
independent differential oracles for the reference implementation: a defect
in its hand-written verification cannot make a non-convergent machine pass
either checker.  The proof scripts are available in the \texttt{coq/}
directory of the accompanying artifact; \texttt{coq/PAPER-MAP.md} maps every
numbered result of this paper to its Coq statement.
\end{remark}

\begin{theorem}[Stream Convergence]
\label{thm:convergence}
Let $R$ be a registry satisfying WFC (\Cref{ax:wfc}) and CC
(\Cref{ax:cc}), and let $P_1, P_2$ be stream processors
(\Cref{def:processor}) consuming the same event stream~$\mathcal{S}$.
Then:
\begin{enumerate}[leftmargin=*,nosep,label=(\alph*)]
\item \textbf{Validity}: For each processor, after it completes the
  incremental compensation phase for any applied event, its state satisfies
  $V_R(\Psi_{P_j}(t)) = \top$.
\item \textbf{Order independence}: Once a processor has finished processing
  its received event set, its state depends only on that set, not on the
  order in which events were applied.
\item \textbf{Agreement}: Processors that have finished processing the same
  events are in the same state: if $P_1$ has finished processing
  $E_{P_1}(t_1)$, $P_2$ has finished processing $E_{P_2}(t_2)$, and
  $E_{P_1}(t_1) = E_{P_2}(t_2)$, then $\Psi_{P_1}(t_1) = \Psi_{P_2}(t_2)$.
  The times $t_1$ and $t_2$ may differ.
\end{enumerate}
\end{theorem}

\begin{proof}
We prove each part.

\medskip\noindent\textbf{Part (a): Validity.}\;
By \Cref{lem:termination}, compensation terminates.  By
\Cref{def:processor}, the processor compensates after each event
application until validity is restored.  Therefore after completing the
compensation phase for any applied event,
$V_R(\Psi_{P_j}(t)) = \top$.

\medskip\noindent\textbf{Part (b): Order independence.}\;
Each processor's computation on event set $E = E_{P_j}(t)$ is a reduction
sequence in the governance rewrite system from $(\sigma_0, E)$; when the
processor has finished, it has reached a normal form $(\sigma^*, \varnothing)$.
Different causality-respecting orders correspond to different reduction
sequences.  By \Cref{cor:unique-nf}, all reduction sequences reach the same
normal form.  Therefore the final state~$\sigma^*$ depends only on $E$, not on
the order.

\medskip\noindent\textbf{Part (c): Agreement.}\;
If $E_{P_1}(t_1) = E_{P_2}(t_2) = E$ and both processors have finished, then
both configurations are normal forms of reductions from $(\sigma_0, E)$.  By
\Cref{cor:unique-nf}, the normal form is unique, so
$\Psi_{P_1}(t_1) = \Psi_{P_2}(t_2)$.
\end{proof}

\begin{remark}[Why ``finished'' is needed]
\label{rem:finished}
Without the qualifier, part~(c) fails: two processors can hold the same
received set while one of them has not yet applied it.  With a counter
registry and the stream $\{e\}$, one processor that has applied~$e$ (state~1)
and one that has received but not yet applied it (state~0, a zero-length
reduction from $(0, \{e\})$) are both processors, both fair, and disagree.
Eventual delivery alone also does not make received sets coincide on an
infinite stream: if $P_1$ receives the $n$-th event at time~$n$ and $P_2$ at
time~$n+1$, both can be finished at every time and still disagree at every
time.  What does hold is (c) as stated, at the same or at different times,
and fixed-point agreement for finite streams (\Cref{cor:quiescent}).  Both
counterexamples are mechanized
(\coq{base_thm_convergence_c_counterexample},
\coq{base_thm_convergence_transient_counterexample}), as is the theorem
(\coq{base_thm_convergence}, \coq{stream_agreement}).
\end{remark}

\begin{remark}[Semantic Projection]
\label{rem:set-function}
The theorem implies that the mapping $F(E) = \nf(\sigma_0, E)$---the unique
normal form of the governance rewrite system starting from $(\sigma_0, E)$---is
a well-defined function of the event \emph{set}~$E$, independent of causal
linearization.  This is the Church--Rosser property for our operational
semantics.
\end{remark}

\begin{corollary}[Convergence Under Quiescent Streams]
\label{cor:quiescent}
If the stream $\mathcal{S}$ is finite and each processor is fair (it keeps
processing while no new events arrive), then there exists a valid state
$\sigma^*$ such that both processors converge to~$\sigma^*$: there exists $T$
with $\Psi_{P_1}(t) = \Psi_{P_2}(t) = \sigma^*$ for all $t \geq T$.
\end{corollary}
\begin{proof}
By eventual delivery, there exists $T_N$ such that
$E_{P_1}(t) = E_{P_2}(t) = \mathcal{S}$ for all $t \geq T_N$.  From $T_N$ on no
new events arrive, so by fairness and \Cref{lem:termination} each processor
finishes processing~$\mathcal{S}$ at some time and stays at that normal form
afterwards.  By \Cref{thm:convergence}(c), from the later of the two finishing
times~$T$ both processors are in the same valid state.
\end{proof}

\begin{remark}[Infinite Streams]
\label{rem:infinite-streams}
For infinite streams, the received event sets grow monotonically
toward~$\mathcal{S}$ but may differ at every time, and eventual delivery does
not make them coincide: a processor that is always one event behind another
disagrees with it forever (\Cref{rem:finished}).  The guarantee is
\Cref{thm:convergence}(c): whenever two finished processors have received the
same events, at the same or at different times, they hold the same state.
The agreed-upon state may continue to change as new events arrive.
\Cref{cor:quiescent} recovers fixed-point convergence when the stream is
finite.
\end{remark}

%========================================
\section{Necessity of the Structural Conditions}
\label{sec:necessity}

We show that UBC is necessary for uniform complexity, and that CC is the exact
condition for agreement where an exact condition of this form exists.

\subsection{Necessity of UBC}

UBC is not merely convenient; it is necessary for a uniform bound on
compensation depth, and thus for the constant per-event compensation overhead
used in the complexity analysis (\Cref{sec:complexity}).  Without it,
compensation depth can grow without limit as event sets grow.

\begin{theorem}[Unbounded Compensation Depth Without UBC]
\label{thm:necessity}
For any $M \in \mathbb{N}$, there exists a registry $R_\infty$ and a finite
event set~$E$ such that some reduction sequence from $(\sigma_0, E)$ in the
governance rewrite system contains more than $M$ compensation steps.
Consequently, no uniform bound on compensation depth exists, and UBC fails.
\end{theorem}

\begin{proof}
Define $R_\infty$ with $\Sigma_{R_\infty} = \mathbb{Z}$, the single invariant
$\psi(\sigma) = \top \iff \sigma = 0$, initial state $\sigma_0 = 0$, and
compensation $\rho(\sigma) = \sigma + 1$ for $\sigma < 0$ and
$\rho(\sigma) = \sigma - 1$ for $\sigma > 0$ (so compensation moves
toward~$0$).

For each $n \in \mathbb{N}$, consider the single-event set $E_n = \{e_n\}$
where $\mathrm{apply}(e_n, 0) = -n$.  Under the processor model
(\Cref{def:processor}), the processor applies $e_n$ to reach state $-n$, then
compensates incrementally: $-n \to -(n-1) \to \cdots \to -1 \to 0$.  This
requires exactly $n$ compensation steps.

For any proposed bound~$M$, choose $n = M + 1$.  The compensation episode
after this single event contains $n > M$ steps.  Therefore no finite~$M$
uniformly bounds compensation depth, and UBC fails.

The measure $\Phi(\sigma) = |\sigma|$ strictly decreases under~$\rho$, so WFC
holds and compensation terminates for each fixed~$E_n$.  The pathology is not
non-termination but \emph{unbounded} termination: the compensation depth
required to restore validity grows without limit.
\end{proof}

\begin{remark}[What Fails]
The counterexample satisfies WFC (compensation terminates for every
fixed~$E_n$) but violates UBC (compensation depth is~$n$, unbounded).  In
fact every reduction from $(0, \{e_n\})$ to a normal form has exactly $n$
compensation steps, and every measure witnessing WFC for $R_\infty$ is
unbounded, not only $|\sigma|$.  CC also holds: the only valid state is~$0$,
so $\rho^*(\sigma) = 0$ for every~$\sigma$, and both sides of CC1 and of CC2
equal~$0$ for every state and every pair of events, whatever
$\mathrm{apply}(e_n, -)$ does away from~$0$ (take, for instance,
$\mathrm{apply}(e_n, \sigma) = \sigma - n$).  The pathology is purely one of
unbounded depth, which invalidates the uniform step bound and the constant
per-event compensation cost in \Cref{thm:complexity}.  Finite semantic state
spaces prevent this by bounding~$\Phi$ (\Cref{lem:finite-implies-ubc}).
\end{remark}

\subsection{Necessity of CC}

Without CC, processors applying the same events in different
causality-respecting orders can reach different valid states
(\Cref{prop:cc-necessary}), and under free delivery this is the only way
agreement can fail (\Cref{thm:cc-exact}).

\begin{proposition}[Disagreement Without CC]
\label{prop:cc-necessary}
There exists a registry satisfying WFC and UBC, and a finite event set~$E$
with two causality-respecting orderings, such that the governance rewrite
system reaches different valid normal forms under the two orderings.
\end{proposition}

\begin{proof}
Let $\Sigma_R = \{A, B, C, D\}$ with valid states $\{A, B\}$ and
compensation $\rho_R(C) = A$, $\rho_R(D) = B$.  UBC holds with $M = 1$.

Define two events $e_1, e_2$ with $\mathrm{deps}(e_1) = \mathrm{deps}(e_2) =
\varnothing$, so both are enabled from the initial configuration
$(A, \{e_1, e_2\})$ and are causally independent under
\Cref{def:causal-order}.  Their application is:
\[
\begin{array}{c|cccc}
  & A & B & C & D \\ \hline
e_1 & C & C & C & D \\
e_2 & D & C & C & D
\end{array}
\]

Evaluate CC1 at $\sigma = A$.  Under the processor model
(\Cref{def:processor}), each processor normalizes after each application,
so the comparison is between $\rho_R^*(A_{e_2}(\rho_R^*(A_{e_1}(A))))$ and
the symmetric order.

\emph{Path 1} ($e_1$ then $e_2$): $\mathrm{apply}(e_1, A) = C$,
$\rho_R^*(C) = A$, $\mathrm{apply}(e_2, A) = D$, $\rho_R^*(D) = B$.
Normal form: $B$.

\emph{Path 2} ($e_2$ then $e_1$): $\mathrm{apply}(e_2, A) = D$,
$\rho_R^*(D) = B$, $\mathrm{apply}(e_1, B) = C$, $\rho_R^*(C) = A$.
Normal form: $A$.

Since $A \neq B$, CC1 is violated.  Two processors consuming
$\{e_1, e_2\}$ in different causality-respecting orders reach different valid
states despite WFC and UBC holding.
\end{proof}

\begin{remark}
The key feature is that $\rho_R$ maps distinct invalid states to
\emph{distinct} valid states ($C \mapsto A$, $D \mapsto B$), creating
path-dependent repair.  Contrast this with the order fulfillment example
(\Cref{sec:example}), where compensation preserves approval intent via the
held state, ensuring path-independent repair.  \Cref{thm:cc-exact} makes
precise in what sense CC captures the boundary between these cases.
\end{remark}

\Cref{prop:cc-necessary} is existential: it shows that \emph{some} failure of
CC produces disagreement.  The converse holds in the following exact form.
Say that \emph{every buffered event may fire} (free delivery) when an event is
enabled as soon as it is in the buffer, as for events without causal
dependencies, and that compensation is \emph{canonical} when $\rho_R^*$
returns a valid state, as \Cref{def:rhostar} guarantees under WFC.  A state is
\emph{reachable} from~$\sigma_0$ if some reduction from some
$(\sigma_0, E)$ passes through it.

\begin{theorem}[CC Is Exact]
\label{thm:cc-exact}
Let $R$ satisfy WFC, with $\rho_R^*$ as in \Cref{def:rhostar}, and let every
buffered event fire freely.  Then every configuration $(\sigma_0, E)$ with $E$
finite has a unique normal form if and only if CC1 and CC2 hold at every
state reachable from~$\sigma_0$.  Under an arbitrary enabledness predicate
(causal or guarded), the system is confluent from a configuration $c_0$ if and
only if, at every configuration reachable from~$c_0$, the governed successors
of every critical pair are joinable; CC1 and CC2 are the special case in which
the join is an equality.
\end{theorem}

\begin{proof}[Proof sketch]
Sufficiency is \Cref{cor:unique-nf}, with Newman's Lemma localized to the
step-closed set of reachable configurations.  For necessity under free
delivery, the two sides of CC1 at~$\sigma$ are the normal forms of the runs
$e_1$ then $e_2$ and $e_2$ then $e_1$ from $(\sigma, \{e_1, e_2\})$, and the
two sides of CC2 at an invalid~$\sigma$ are the normal forms of
apply-then-compensate and compensate-then-apply from $(\sigma, \{e\})$.  A
reachable~$\sigma$ is reached from $\sigma_0$ by delivering some word~$w$
first, so a failure of CC at~$\sigma$ yields two distinct normal forms of
$(\sigma_0, w \cdot e_1 e_2)$ or $(\sigma_0, w \cdot e)$.  The general case is
Newman's Lemma together with the observation that a non-joinable critical pair
at a reachable configuration is a pair of distinct normal forms.  Both
directions are mechanized (\coq{cc_exact_from}, \coq{wfc_cc_exact_from},
\coq{jc_exact}).
\end{proof}

\begin{remark}[Scope of the Converse]
\label{rem:converse-scope}
Both qualifiers are needed.  If $\rho_R^*$ is only required to be some
compensation reduct of its argument (the identity qualifies), CC1 can fail
while every configuration has a unique normal form.  And under causal
enabledness, CC1 can fail for a co-enabled pair at a reachable configuration
while every run still reaches the same normal form: events $M_1$ (write~1)
and $M_2$ (write~2) do not commute, but a reset event that depends on both
erases the difference, so from $(0, \{M_1, M_2, M_{\mathrm{reset}}\})$ every
normal form is $(0, \varnothing)$ (\coq{rho_star_qualifier},
\coq{masked_cc1}).  There the exact condition is the
joinability form of \Cref{thm:cc-exact}, not CC1.  Together with
\Cref{thm:necessity}, which satisfies CC and violates UBC, and
\Cref{prop:cc-necessary}, which satisfies WFC and UBC and violates CC, this
shows each condition is needed for its guarantee: UBC for uniform step
bounds, CC (exactly, under free delivery) for agreement.
\end{remark}

%========================================
\section{Complexity Analysis}
\label{sec:complexity}

\begin{theorem}[Convergence Complexity]
\label{thm:complexity}
For a stream prefix of length $n$ under a registry satisfying UBC with
bound~$M$ and $k = |I_R|$ invariants, the model-level cost of
registry-governed stream processing is $O(n)$ for event application,
validation, and compensation combined.  If causality maintenance is
implemented with a priority queue, the total cost is $O(n \log n)$, where $k$
and $M$ are treated as constants of system design.
\end{theorem}

\begin{proof}
We analyze the per-event cost, separating model-level costs (which follow from
the formalism) from scheduling costs (which depend on implementation).

\textbf{Model-level per-event costs.}  For each event $e_i$
($1 \leq i \leq n$):
\begin{itemize}[leftmargin=*,nosep]
\item \emph{Event application}: $O(1)$ (state update).
\item \emph{Invariant validation}: $O(k) = O(1)$ (evaluate each~$\psi_j$).
\item \emph{Compensation}: At most $M$ steps, each $O(1)$.  Since $M$ is
  bounded by UBC, this is $O(1)$.
\end{itemize}
Registry governance therefore adds $O(1)$ overhead per event beyond what the
base stream processor already performs.

\textbf{Scheduling overhead.}  Maintaining a set of enabled events and
resolving causal dependencies requires implementation-dependent data
structures.  With a priority queue keyed by timestamp and identifier, each
event insertion and dependency-resolution check costs $O(\log i)$, yielding:

\textbf{Total cost.}
\[
T(n) = \sum_{i=1}^n \bigl(\underbrace{O(1)}_{\text{model}} +
  \underbrace{O(\log i)}_{\text{scheduling}}\bigr)
     = O(n) + O\!\left(\sum_{i=1}^n \log i\right)
     = O(n \log n).
\]

The dominant term is causality maintenance.  Validation and compensation
contribute no asymptotic overhead.  Systems with native causal ordering (e.g.,
single-partition Kafka topics) eliminate the $O(\log i)$ term, reducing the
bound to~$O(n)$.
\end{proof}

%========================================
\section{Discussion}
\label{sec:discussion}

\subsection{Three Regimes of Coordination-Free Convergence}
\label{sec:design-space}

Our result, combined with CRDTs~\cite{shapiro2011crdt} and invariant
confluence~\cite{bailis2014coordination}, identifies three distinct structural
regimes under which distributed processors converge without coordination:

\begin{enumerate}[leftmargin=*,nosep]
\item \textbf{Operation commutativity} (CRDTs): Operations commute
  algebraically.  Convergence follows from the algebra of operations.  No
  invariant enforcement or compensation is needed.
\item \textbf{Invariant confluence} (I-confluence): Operations preserve
  application invariants under all orderings.  Convergence follows because
  validity is never violated.  No compensation is needed.
\item \textbf{Normalization confluence} (this paper): Operations may be
  non-commutative and may violate invariants.  Compensation normalizes states
  to validity.  WFC and CC are sufficient for semantic convergence
  (\Cref{thm:convergence}); UBC is orthogonal to convergence and governs
  scalability (\Cref{thm:complexity}).  No coordination is needed.
\end{enumerate}

The regimes relate as follows.  Operation commutativity with every state
valid implies CC (with $\rho_R^* = \mathrm{id}$), and conversely a
compensation-free system meets the causal convergence condition exactly when
it is an op-based CRDT, so normalization confluence strictly contains the
op-based CRDTs.  Invariant confluence implies that compensation never fires,
but it does not by itself make operations commute: in the
sequential-application model of this paper, an invariant-confluent system
satisfies CC exactly when its operations commute on the reachable states, and
by \Cref{thm:cc-exact} one whose operations do not commute there diverges
under free delivery.  So normalization confluence contains the
invariant-confluent systems that converge under reordering, not all of them,
and it additionally covers systems where operations are non-commutative and
invariant-violating, the common case in business applications.

The regimes partition the coordination-free design space by \emph{where}
convergence is guaranteed: in the operations (regime~1), in the invariants
(regime~2), or in the normalization rewrite system (regime~3).

\begin{remark}[Convergence as Rewrite Property]
\label{rem:rewrite-property}
The conventional literature frames convergence as a property of operations:
operations must commute, or operations must preserve invariants.
Normalization confluence reframes convergence as a property of the
\emph{rewrite system} that interleaves application and compensation.  The
operations themselves need no special algebraic structure.  This shifts the
design burden from operation algebra to compensation design, and places the
convergence mechanism at the semantic normalization layer---above the
data-structure layer (CRDTs) and the transaction layer (I-confluence).
\end{remark}

\subsection{Verifying Compensation Commutativity}
\label{sec:verify-cc}

CC is easy to state but potentially hard to verify: it quantifies over all
states and all independent event pairs.  \Cref{sec:calculus} develops a
verification calculus that reduces CC to local per-event and per-pair checks
via invariant footprints and decomposable repair.  Here we note three common
patterns that practitioners can use as a first check before invoking the full
calculus:

\begin{enumerate}[leftmargin=*,nosep]
\item \textbf{Compensation to a canonical state.}\;  Suppose there is a
  default valid state $\sigma_\bot$ such that $\rho_R^*(\sigma) = \sigma_\bot$
  for all invalid~$\sigma$ (the ``reset to safe default'' pattern).  This
  alone does not give CC: on valid states compensation never fires, so two
  valid-to-valid events that do not commute violate CC1, and an event that
  treats an invalid state differently from the reset state violates CC2.
  What the pattern does give is a cheap check.  CC2 holds for~$e$ if and only
  if $N(A_e(\sigma)) = N(A_e(\sigma_\bot))$ for every invalid~$\sigma$ (one
  comparison per invalid state against a fixed reference), where
  $N = \rho_R^*$ and $A_e = \mathrm{apply}(e, -)$ as in
  \Cref{sec:calculus}; and when this holds for $e_1$ and~$e_2$, CC1 for the
  pair at~$\sigma$ holds if and only if
  $N(A_{e_2}(A_{e_1}(\sigma))) = N(A_{e_1}(A_{e_2}(\sigma)))$.  So canonical
  compensation, the reset check, and commutation up to~$N$ together give CC1
  and CC2.
\item \textbf{Lattice-structured compensation.}\;  Suppose the states form a
  join-semilattice, each event joins a fixed element
  ($\mathrm{apply}(e, \sigma) = \sigma \sqcup a_e$), the normalizer~$N$ is a
  closure operator (monotone, extensive, idempotent), and one compensation
  step stays between a state and its closure
  ($\sigma \sqsubseteq \rho_R(\sigma) \sqsubseteq N(\sigma)$).  Then
  $N(N(x) \sqcup y) = N(x \sqcup y)$ for all $x, y$, so both sides of CC1
  equal $N(\sigma \sqcup a_{e_1} \sqcup a_{e_2})$, and CC2 holds because
  $\sigma \sqcup a_e \sqsubseteq \rho_R(\sigma) \sqcup a_e \sqsubseteq
  N(\sigma) \sqcup a_e$.  (This is a short paper-level argument; it is not
  part of the mechanization.)  Monotone repair on an ordered shared domain is
  also the setting of the companion paper's~\cite{blackwell2026federated}
  cyclic federations, where federated repair reaches a least fixed point (by
  the Knaster--Tarski theorem) and a further condition decides whether event
  order matters, with state-based CRDTs as the compensation-free special case.
\item \textbf{Decomposable repair.}\;  If compensation decomposes into
  commuting per-invariant repairs, CC can be verified via the footprint
  calculus of \Cref{sec:calculus}.  The order fulfillment example
  (\Cref{sec:example}) uses this pattern, with one invariant whose repair is
  the hold.
\end{enumerate}

\subsection{Applications}

The compensate-and-converge principle applies wherever systems process streams
of changes under semantic constraints with bounded compensation.  We highlight
three domains.

\paragraph{Infrastructure Configuration Drift.}
Infrastructure-as-code tools apply streams of configuration changes to
infrastructure state.  Valid configurations are defined by policy.  The
semantic state space is finite ($k$ boolean policy predicates yield at most
$2^k$ equivalence classes).  Remediation operators that reset violated policies
to compliant defaults fit the canonical-state pattern: CC2 reduces to one
comparison per violating state against the default, and CC1 to commutation of
changes up to normalization.  When those checks pass, our theorem guarantees
convergence to equivalent policy-compliant states regardless of the order in
which changes are applied.

\paragraph{Database Schema Migration.}
Schema migrations are event streams applied to schema state.  The set of valid
schema configurations is finite.  Migration scripts that repair constraint
violations constitute compensation operators.  When such scripts resolve each
constraint independently (decomposable repair), CC reduces to the per-event
and per-pair checks of the calculus (\Cref{sec:calculus}); when those pass,
our theorem guarantees order-independent convergence.

\paragraph{Regulatory Compliance.}
Organizations processing streams of regulatory changes must maintain
compliance.  Compliance states are finite (each regulation is either satisfied
or not).  Remediation procedures that restore compliance by addressing each
violated regulation independently are decomposable repairs, so CC can be
checked with the calculus; when it holds, our theorem guarantees convergence
to equivalent compliance states.

\subsection{Infinite Domains}
\label{sec:infinite}

Nothing in the convergence result requires $\Sigma_R$ to be finite.  Termination
(\Cref{lem:termination}) rests on the well-founded measure $\mu = (|B|, \Phi)$:
the buffer $B$ is finite because the received event set is, and $\Phi$ maps into
a well-founded order---neither depends on $|\Sigma_R| < \infty$.  Local
confluence (\Cref{lem:local-confluence}) requires CC to \emph{hold}, not to be
enumerable.  Convergence is therefore domain-independent.

\begin{corollary}[Domain-Independent Convergence]
\label{cor:infinite}
\Cref{thm:convergence} and \Cref{cor:quiescent} hold for any semantic state
space, finite or infinite, provided CC holds and WFC holds with a measure
$\Phi$ into any well-founded order (not necessarily~$\mathbb{N}$).
\end{corollary}

This is mechanized at both the rewrite-system and the stream level, with an
instance on~$\mathbb{Z}$ (\Cref{sec:mechanization}).  The necessity example of
\Cref{sec:necessity} also witnesses it: the registry $R_\infty$
($\Sigma = \mathbb{Z}$, $\Phi(\sigma) = |\sigma|$, compensation stepping
toward~$0$) satisfies WFC and CC and converges (every event set reaches a
unique valid normal form), yet $\Sigma$ is infinite.  What it lacks is UBC:
compensation depth grows without bound.

Finiteness therefore underwrites three \emph{secondary} properties, not
soundness: (i)~UBC and the uniform per-event cost (\Cref{thm:complexity});
(ii)~verification by exhaustive enumeration; and (iii)~the precomputed
normal-form and step tables behind $O(1)$ runtime application.  A registry whose
logic is captured by $k$ boolean predicates already has a finite semantic
quotient---$|\Sigma_R| \le 2^k$ (\Cref{rem:scope})---so ``infinite'' here means
genuinely \emph{unbounded numeric} semantic state, as in $R_\infty$.

On such a domain enumeration is unavailable, but CC can still be discharged.
The verification calculus (\Cref{sec:calculus}), with decomposable repair,
strong absorption, footprint absorption and commutation up to normalization,
establishes CC by \emph{algebraic} conditions that never quantify over an
enumerated state space; symbolic decision procedures apply where those do
not.  When the valid states form a complete lattice and compensation is
monotone, the normal form is a fixed point computed by Kleene iteration:
finitely under the ascending chain condition, and under
\emph{widening}~\cite{cousot1977} otherwise, at the cost of a sound
over-approximation.  Without ACC the least fixed point still exists (by the
Knaster--Tarski theorem) but finite iteration need not reach it; whenever WFC
and CC hold, existence and uniqueness of the normal form are guaranteed by
\Cref{cor:infinite}.
The federated generalization of the monotone case
appears in~\cite{blackwell2026federated}.

\subsection{Extensions}

The single-registry result is the foundation for further extensions.  First,
\emph{federated streams} where multiple organizations maintain independent
registries connected by morphisms; this is developed in the companion
paper~\cite{blackwell2026federated}, which gives exact conditions for
federated convergence when events fire at federally valid states, and treats
multi-source networks (resolution operators) and cyclic networks under
monotone repair.
Second, \emph{privacy-preserving validation} where organizations prove registry
compliance via zero-knowledge proofs, which remains open.  Both build on the
single-registry convergence result established here.

%========================================
\section{Verification Calculus for CC}
\label{sec:calculus}

Compensation commutativity is a global condition: it quantifies over all states
and all independent event pairs.  This section develops conditions, exact
where possible, that reduce CC to local per-event and per-pair checks.  The key tools are an
\emph{invariant footprint} mapping and a \emph{decomposable repair} condition.

\subsection{Notation}

For readability in algebraic manipulations, we write
$N(\sigma) \defeq \rho_R^*(\sigma)$ for iterated compensation to validity
(\Cref{def:rhostar}).  The normalizer satisfies: (i)~$V_R(N(\sigma)) = \top$
for all~$\sigma$; (ii)~$N(\sigma) = \sigma$ if $V_R(\sigma) = \top$;
(iii)~$N(N(\sigma)) = N(\sigma)$ (idempotence) (\coq{rho_star_valid},
\coq{rho_star_fix}, \coq{rho_star_idem}).  We write
$A_e(\sigma) \defeq \mathrm{apply}(e, \sigma)$ when convenient.

In this notation, the CC obligations become:
\begin{itemize}[leftmargin=*,nosep]
\item[\textbf{(CC1)}] $N(A_{e_2}(N(A_{e_1}(\sigma)))) =
  N(A_{e_1}(N(A_{e_2}(\sigma))))$ for independent $e_1, e_2$.
\item[\textbf{(CC2)}] $N(A_e(\sigma)) = N(A_e(\rho_R(\sigma)))$ for
  invalid~$\sigma$.
\end{itemize}
We say that CC1 holds for $(e_1, e_2)$ \emph{at}~$\sigma$ when its equation
holds at that state, and that CC2 holds \emph{for}~$e$ when its equation holds
at every invalid~$\sigma$.

\subsection{Invariant Vectors and Footprints}

\begin{definition}[Invariant Vector]
\label{def:inv-vector}
Define $\mathbf{v}: \Sigma_R \to \{0,1\}^k$ by
$\mathbf{v}(\sigma)_i = 1 \iff \psi_i(\sigma) = \top$.
\end{definition}

\begin{definition}[Footprint]
\label{def:footprint}
A \emph{footprint} is a function $F: \mathcal{E} \to 2^{\{1,\ldots,k\}}$ such
that for all $\sigma$ and all $i \notin F(e)$,
\[
\mathbf{v}(A_e(\sigma))_i = \mathbf{v}(\sigma)_i.
\]
Intuitively, $F(e)$ is the set of invariants whose truth value can change
under~$e$.  Over-approximations are sound: larger footprints make the calculus
more conservative but never incorrect.
\end{definition}

\subsection{Decomposable Repair}

We capture the common case where repair addresses each violated invariant
independently.

\begin{definition}[Per-Invariant Repair]
\label{def:perinv}
A set of \emph{per-invariant repair operators} $r_1, \ldots, r_k :
\Sigma_R \to \Sigma_R$ satisfies:
\begin{enumerate}[leftmargin=*,nosep,label=(R\arabic*)]
\item \textbf{Targeted fix:}\; if $\psi_i(\sigma) = \bot$, then
  $\psi_i(r_i(\sigma)) = \top$.
\item \textbf{No-op on satisfied:}\; if $\psi_i(\sigma) = \top$, then
  $r_i(\sigma) = \sigma$.
\item \textbf{Mutual commutativity:}\; $r_i(r_j(\sigma)) = r_j(r_i(\sigma))$
  for all $i, j$.
\end{enumerate}
\end{definition}

\begin{definition}[Decomposable Normalizer]
\label{def:decomp}
The normalizer $N$ is \emph{decomposable} if there exist per-invariant repairs
$r_1, \ldots, r_k$ as above such that
\[
N(\sigma) = r_1(r_2(\cdots r_k(\sigma) \cdots))
\]
for all~$\sigma$.  (The order is irrelevant by mutual commutativity.)
\end{definition}

\begin{definition}[Repair Locality]
\label{def:repair-locality}
An event $e$ is \emph{repair-local} with respect to footprint $F(e)$ if for
all $i \notin F(e)$,
\[
r_i(A_e(\sigma)) = A_e(r_i(\sigma)).
\]
That is, $e$ commutes with repairs of invariants it does not affect.
\end{definition}

\subsection{Strong Absorption and CC2}

\begin{theorem}[Strong Absorption]
\label{thm:strong-absorption}
Assume WFC and let $N = \rho_R^*$.  Then CC2 holds (for all events) if and
only if, for all events $e$ and states~$\sigma$,
\[
N(A_e(\sigma)) = N(A_e(N(\sigma))).
\]
\end{theorem}

\begin{proof}
($\Leftarrow$) Let $V_R(\sigma) = \bot$.  Because $N$ iterates $\rho_R$ to the
unique valid fixpoint reachable from~$\sigma$, and $\rho_R(\sigma)$ lies on the
same compensation chain, we have $N(\sigma) = N(\rho_R(\sigma))$.  Applying the
assumed equation at $\sigma$ and at $\rho_R(\sigma)$:
\[
N(A_e(\sigma)) = N(A_e(N(\sigma))) = N(A_e(N(\rho_R(\sigma)))) = N(A_e(\rho_R(\sigma))),
\]
where the first and third equalities use the assumption and the middle equality
uses $N(\sigma) = N(\rho_R(\sigma))$.  This is CC2.

($\Rightarrow$) This is the generalization of CC2 noted in
\Cref{rem:absorption}: by induction on the number~$m$ of compensation steps
from $\sigma$ to $N(\sigma) = \rho_R^m(\sigma)$, CC2 at each invalid
$\rho_R^j(\sigma)$, $j < m$, gives
$N(A_e(\rho_R^j(\sigma))) = N(A_e(\rho_R^{j+1}(\sigma)))$, and chaining these
equalities gives $N(A_e(\sigma)) = N(A_e(N(\sigma)))$.
\end{proof}

Both directions are mechanized (\coq{strong_absorption_iff_cc2};
the forward direction alone is \coq{base_thm_strong_absorption}).  The
per-event form also holds: CC2 for a single event~$e$ at every invalid state
implies strong absorption for~$e$ (\coq{calc_cc2_strong_absorption}).

\begin{remark}
Strong absorption states: ``normalizing before applying an event does not
change the final normalized outcome.''  This holds whenever compensation
repairs violations in a way determined by the violation pattern rather than by
the path to the violation.  The order fulfillment example
(\Cref{sec:example}) satisfies this: the held state captures the approval
intent regardless of whether the order was directly approved with low balance
or arrived at the invalid state through a different path.
\end{remark}

\subsection{Footprints and CC1}

We prove two preparatory lemmas before the main theorem.

\begin{lemma}[Repair-Event Commutation]
\label{lem:repair-commute}
If event $e$ is repair-local with respect to footprint $F(e)$
(\Cref{def:repair-locality}), then for any index set $J$ with
$J \cap F(e) = \varnothing$ and any state~$\sigma$:
\[
R_J(A_e(\sigma)) = A_e(R_J(\sigma)),
\]
where $R_J = \prod_{i \in J} r_i$.
\end{lemma}

\begin{proof}
By repair locality, $r_i \circ A_e = A_e \circ r_i$ for each $i \in J$
(since $i \notin F(e)$).  The result follows by induction on $|J|$: if the
claim holds for $R_{J'}$ with $|J'| < |J|$, then for any $j \in J$,
\[
R_J(A_e(\sigma)) = r_j(R_{J \setminus \{j\}}(A_e(\sigma)))
= r_j(A_e(R_{J \setminus \{j\}}(\sigma)))
= A_e(r_j(R_{J \setminus \{j\}}(\sigma)))
= A_e(R_J(\sigma)),
\]
using the inductive hypothesis and then $r_j \circ A_e = A_e \circ r_j$.
\end{proof}

\begin{lemma}[Repair Block Idempotence]
\label{lem:repair-idempotent}
For any index set $J \subseteq \{1,\ldots,k\}$, the repair block $R_J$ is
idempotent on states that are valid on all invariants in~$J$: if
$\psi_i(\sigma) = \top$ for all $i \in J$, then $R_J(\sigma) = \sigma$.
More generally, $R_J(R_J(\sigma)) = R_J(\sigma)$ for all~$\sigma$.
\end{lemma}

\begin{proof}
If $\psi_i(\sigma) = \top$ for all $i \in J$, then each $r_i(\sigma) = \sigma$
by~(R2), so $R_J(\sigma) = \sigma$.  For general idempotence: $R_J(\sigma)$
satisfies $\psi_i(R_J(\sigma)) = \top$ for all $i \in J$ (each $r_i$ fixes its
invariant by~(R1) and does not break others by mutual commutativity and~(R2)),
so $R_J(R_J(\sigma)) = R_J(\sigma)$ by the first claim.
\end{proof}

\begin{theorem}[Footprint Commutativity]
\label{thm:footprint-cc1}
Let $N = \rho_R^*$ be decomposable via per-invariant repairs
$r_1, \ldots, r_k$ satisfying (R1) to (R3), let $F$ be a footprint, and let
$e_1, e_2$ be events.  Write $F_j = F(e_j)$ and consider the hypotheses
\begin{enumerate}[leftmargin=*,nosep,label=(H\arabic*)]
\item Disjoint footprints: $F_1 \cap F_2 = \varnothing$.
\item Repair locality: each $e_j$ is repair-local with respect to $F_j$.
\item Commutation up to~$N$ at~$\sigma$:
  $N(A_{e_2}(A_{e_1}(\sigma))) = N(A_{e_1}(A_{e_2}(\sigma)))$ (raw
  commutation $A_{e_1} \circ A_{e_2} = A_{e_2} \circ A_{e_1}$ suffices).
\item Footprint absorption: $N(A_{e_j}(R_{F_j}(\tau))) = N(A_{e_j}(\tau))$
  for all states~$\tau$ and $j = 1, 2$.
\end{enumerate}
Then:
\begin{enumerate}[leftmargin=*,nosep,label=(\alph*)]
\item \textbf{Exact form.}  Under (H1) and (H2), for every state~$\sigma$,
  CC1 holds at~$\sigma$ if and only if
  $N(A_{e_2}(A_{e_1}(R_{F_2}(\sigma)))) = N(A_{e_1}(A_{e_2}(R_{F_1}(\sigma))))$.
\item \textbf{Valid states.}  Under (H1) and (H2), for every valid~$\sigma$,
  CC1 holds at~$\sigma$ if and only if (H3) holds at~$\sigma$.  In
  particular, (H1), (H2) and raw commutation give CC1 at every valid state.
\item \textbf{Every state.}  Under (H2), (H3) at~$\sigma$ and (H4), CC1 holds
  at~$\sigma$; (H1) is not needed.  Given (H2), footprint absorption (H4)
  for~$e_j$ is equivalent to strong absorption for~$e_j$:
  $N(A_{e_j}(\tau)) = N(A_{e_j}(N(\tau)))$ for all~$\tau$.
\end{enumerate}
\end{theorem}

\begin{proof}
Write $J^c$ for $\{1, \ldots, k\} \setminus J$.  First, $N \circ R_J = N$ for
every index set~$J$: by decomposability and (R3), $N = R_{J^c} \circ R_J$, so
$N(R_J(\tau)) = R_{J^c}(R_J(R_J(\tau))) = R_{J^c}(R_J(\tau)) = N(\tau)$ by
\Cref{lem:repair-idempotent}.  Second, if $e$ is repair-local with respect to
$F(e)$, then for every~$\tau$
\begin{equation}
\label{eq:nan}
N(A_e(N(\tau))) = N(A_e(R_{F(e)}(\tau))),
\end{equation}
because $N(\tau) = R_{F(e)^c}(R_{F(e)}(\tau))$, \Cref{lem:repair-commute}
moves $R_{F(e)^c}$ past $A_e$, and $N \circ R_{F(e)^c} = N$.

(a) By \eqref{eq:nan} for~$e_2$, the left side of CC1 is
$N(A_{e_2}(N(A_{e_1}(\sigma)))) = N(A_{e_2}(R_{F_2}(A_{e_1}(\sigma))))$, and
by (H1) and \Cref{lem:repair-commute} for~$e_1$ (with $J = F_2$, disjoint
from~$F_1$) this is $N(A_{e_2}(A_{e_1}(R_{F_2}(\sigma))))$.  Symmetrically,
the right side is $N(A_{e_1}(A_{e_2}(R_{F_1}(\sigma))))$.

(b) At a valid~$\sigma$ every $\psi_i(\sigma)$ holds, so
$R_{F_1}(\sigma) = R_{F_2}(\sigma) = \sigma$ by
\Cref{lem:repair-idempotent}, and (a) becomes (H3).

(c) By \eqref{eq:nan}, (H4) for~$e_j$ says exactly
$N(A_{e_j}(N(\tau))) = N(A_{e_j}(\tau))$ for all~$\tau$, which is strong
absorption (conversely, strong absorption at $R_{F_j}(\tau)$ together with
$N(R_{F_j}(\tau)) = N(\tau)$ gives (H4)).  Then
\[
N(A_{e_2}(N(A_{e_1}(\sigma)))) = N(A_{e_2}(A_{e_1}(\sigma)))
= N(A_{e_1}(A_{e_2}(\sigma))) = N(A_{e_1}(N(A_{e_2}(\sigma)))),
\]
by strong absorption for~$e_2$, (H3), and strong absorption for~$e_1$.
\end{proof}

\begin{corollary}[CC1 Under CC2]
\label{cor:cc1-under-cc2}
Assume WFC and let $N = \rho_R^*$.  If CC2 holds for $e_1$ and~$e_2$, then
for every state~$\sigma$, CC1 holds for the pair at~$\sigma$ if and only if
$N(A_{e_2}(A_{e_1}(\sigma))) = N(A_{e_1}(A_{e_2}(\sigma)))$.
\end{corollary}

\begin{proof}
By \Cref{thm:strong-absorption} (in its per-event form) both events satisfy
strong absorption, and the chain of equalities in the proof of
\Cref{thm:footprint-cc1}(c) applies in both directions.
\end{proof}

In a registry satisfying CC2, then, CC1 is exactly commutation up to
normalization, and footprints play no role in that equivalence; their role is
to make (H4), or CC2 itself, a local check.  Every part is mechanized:
(a)~\coq{calc_footprint_cc1_iff}, (b)~\coq{calc_footprint_cc1_valid_iff}
and \coq{calc_footprint_cc1_valid}, (c)~\coq{calc_footprint_cc1} and
\coq{calc_fp_absorb_iff_sa}, \Cref{cor:cc1-under-cc2}
\coq{calc_cc1_iff_commute_under_cc2}; with non-vacuity instances
\coq{calc_clamp_instance} (two clamped counters with disjoint footprints, CC1
at every state including invalid ones) and \coq{calc_order_cc1}.  Form~(b) is
the one the reference implementation's compositional check uses, since it
checks CC1 on valid states.  The same reasoning covers a state space split
into components: on $\Sigma_1 \times \Sigma_2$ with $N = N_1 \times N_2$ and
events acting on one component each, CC1 holds at every valid state, and at
any state exactly when each event absorbs its own component's normalizer
(\coq{calc_components_cc1_valid}, \coq{calc_components_cc1_iff}).

\begin{remark}[Each Hypothesis Is Needed]
\label{rem:footprint-counterexamples}
Version~1 of this paper stated \Cref{thm:footprint-cc1} with (H1) and (H2)
alone as giving CC1 at every state.  That is false.  With no invariants
($k = 0$) every footprint is empty, repair is the identity, (H1) and (H2) hold
vacuously, and the events $x := 1$ and $x := 2$ violate CC1
(\coq{base_thm_footprint_cc1_refuted}).  Adding raw commutation does not
help: on states $\{0,1,2,3\}$ with $\psi_1 = \psi_2 = (\sigma \geq 2)$,
repairs $r_1 = r_2$ sending $0, 1 \mapsto 2$ and fixing $2, 3$, and events
$a = (0 \mapsto 0, 1 \mapsto 0, 2 \mapsto 2, 3 \mapsto 3)$ (footprint empty,
repair-local) and $b = (0 \mapsto 0, 1 \mapsto 0, 2 \mapsto 3, 3 \mapsto 0)$
(footprint $\{1, 2\}$), the raw events commute but at state~$0$ the two sides
of CC1 are $3$ and~$2$ (\coq{base_thm_footprint_cc1_raw_refuted}).  The
failure is that $b$ does not absorb the repairs of its own footprint
($N(b(R_{\{1,2\}}(0))) = 3 \neq 2 = N(b(0))$), which is (H4); $b$ also
violates CC2 at~$0$.  Dropping any one hypothesis of the corrected forms
breaks them, with every other hypothesis holding: the valid-state form
without (H3) (\coq{calc_valid_needs_commute}), without (H1)
(\coq{calc_valid_needs_disjoint}), or without (H2)
(\coq{calc_valid_needs_repair_local}); the every-state form without (H3)
(\coq{calc_all_needs_commute}), without (H4)
(\coq{calc_all_needs_absorb}, the example above), or without (H2)
(\coq{calc_all_needs_repair_local}).
\end{remark}

\begin{remark}[Practical import]
\label{rem:practical-import}
A practitioner does not prove CC1 by global state enumeration.  Instead:
(i)~assign each event type an invariant footprint; (ii)~verify that each
event commutes with repairs outside its footprint (repair locality,
\Cref{lem:repair-commute}); (iii)~verify, per event, that it absorbs the
repairs of its own footprint, (H4), equivalently strong absorption, which CC2
already implies; and (iv)~verify, per independent pair, commutation up to
normalization, (H3).  All four are local obligations, and (iv) is often raw
commutation.  At valid states, (ii) and disjoint footprints reduce CC1 to
(iv) alone.  The order fulfillment example (\Cref{sec:example}) does not have
disjoint footprints: approval can falsify the single invariant and credit can
restore it, so every footprint of either event contains it
(\coq{base_rem_practical_import_refuted}).  It is certified by part~(c)
instead: the held-state repair commutes with credit application
(\coq{calc_order_credit_repair_commute}), both events absorb their footprint
repairs, and the events commute up to normalization, so CC1 holds at every
state, and CC2 holds as well (\coq{calc_order_cc1}).
\end{remark}

\subsection{Product Composition}

Real systems comprise multiple semantic modules.  The following theorem shows
CC verifications compose.

\begin{definition}[Product Registry]
\label{def:product}
Let $R_1 = (\Sigma_1, I_1, V_1, \rho_1)$ and $R_2 = (\Sigma_2, I_2, V_2,
\rho_2)$ with disjoint invariant sets.  Define $R_1 \otimes R_2$ by:
\[
\Sigma = \Sigma_1 \times \Sigma_2, \quad
V(\sigma_1, \sigma_2) = V_1(\sigma_1) \wedge V_2(\sigma_2), \quad
\rho(\sigma_1, \sigma_2) = (\rho_1(\sigma_1), \rho_2(\sigma_2)).
\]
Events are pairs $e = (e^{(1)}, e^{(2)})$ with application componentwise.
\end{definition}

\begin{theorem}[Product Lifting]
\label{thm:product}
If each $R_j$ satisfies WFC and CC, then $R_1 \otimes R_2$ satisfies WFC and
CC; more precisely, CC1 holds for a product pair at $(\sigma_1, \sigma_2)$ if
and only if it holds for the first components at~$\sigma_1$ and for the
second components at~$\sigma_2$.  If each normalizer is decomposable, the
product normalizer is decomposable with footprints composed by union across
components, and repair locality and footprint disjointness lift.
\end{theorem}

\begin{proof}
\emph{WFC}: Use measure $\Phi(\sigma_1, \sigma_2) = \Phi_1(\sigma_1) +
\Phi_2(\sigma_2)$.  Each compensation step strictly decreases at least one
component when invalid.

\emph{CC2}: The product normalizer satisfies $N(\sigma_1, \sigma_2) =
(N_1(\sigma_1), N_2(\sigma_2))$ and application is componentwise, so
CC2 holds iff it holds in each component.

\emph{CC1}: Both sides of the CC1 equation decompose into paired equalities
in $R_1$ and $R_2$ by componentwise application and normalization.

\emph{Decomposability}: Per-invariant repairs in each component act on
disjoint state components and thus commute across components trivially.
The normalizer identity $N = N_1 \times N_2$ uses that each $\rho_j$ fixes
valid states (\Cref{def:registry}).  Mechanized as \coq{base_thm_product}
and \coq{base_thm_product_decomposable}, with the instance
\coq{calc_product_instance}.
\end{proof}

\begin{remark}
Product composition enables modular system design: verify CC for each semantic
module independently, then compose without re-verification.  An order
processing system with separate inventory, payment, and shipping registries
can verify CC per-module and lift to the combined system.
\end{remark}

\subsection{Partial Decomposability}

Full decomposability is a strong sufficient condition.  When it does not hold
globally, the calculus still provides value:

\begin{enumerate}[leftmargin=*,nosep]
\item The footprint theorem (\Cref{thm:footprint-cc1}) is a statement about
  one event pair: it applies to any pair satisfying its hypotheses,
  regardless of whether \emph{all} pairs do, and \Cref{cor:cc1-under-cc2}
  needs no decomposability at all.
\item When footprints overlap for some event pair, the calculus identifies
  precisely which pairs require further analysis.  A practitioner can then
  check the corrected canonical-state pattern for the overlapping invariants
  (\Cref{sec:verify-cc}, pattern~1) or add
  coordination only for the interacting event types.
\item For event types that share footprints but whose normalized outcomes
  agree by domain-specific reasoning (e.g., both repairs clamp to the same
  bound), CC1 can be verified by case analysis restricted to the overlapping
  states.  The footprint calculus reduces the case space.
\end{enumerate}

%========================================
\section{Conclusion}
\label{sec:conclusion}

We have identified \emph{normalization confluence} as a third structural regime
for coordination-free convergence in distributed systems, alongside operation
commutativity (CRDTs) and invariant confluence (I-confluence).  The regime is
characterized by two conditions: well-founded compensation (WFC), which
ensures termination, and compensation commutativity (CC), which ensures
confluence.  Under both, the governance rewrite system has unique normal forms
by Newman's Lemma, and all processors consuming the same events converge to
the same valid state.  A uniform bound on compensation depth (UBC) additionally
guarantees constant per-event overhead.  Without UBC, compensation depth is
unbounded.  CC is exact: when every buffered event may fire and compensation
repairs to validity, every event set has a unique normal form if and only if
CC holds on the reachable states, and under causal or guarded enabledness a
joinability form of CC is the exact condition.  Convergence does not depend on
the state space being finite; finiteness serves only uniform cost and
checking by enumeration.

To make CC practically verifiable, we developed a verification calculus that
reduces the global CC obligation to local checks.  Invariant footprints
identify which invariants each event can affect; decomposable repair ensures
per-invariant fixes commute; CC2 is equivalent to strong absorption; CC1 at
every state follows from repair locality, per-event footprint absorption and
commutation up to normalization, and at valid states disjoint footprints and
repair locality reduce it to commutation up to normalization alone; and
product composition enables modular verification across semantic modules.
Each hypothesis is shown to be needed.  Together, these convert
``compensate-and-converge'' from a convergence theorem into a verifiable
design pattern with tractable proof obligations.  The convergence theory and
the calculus are mechanized in Coq/Rocq without axioms, including the paper's examples and the
counterexamples that delimit each theorem.

Compensation commutativity is strictly weaker than the operation commutativity
required by CRDTs: operations need not commute, only their compensated results.
The gap between these two conditions is precisely the design space that
normalization confluence occupies---systems where operations are
non-commutative, invariant-violating, but convergent under repair.

The three regimes partition the coordination-free design space by where
convergence is guaranteed: in the operations, in the invariants, or in the
normalization rewrite system.  This last regime formalizes a pattern that
practitioners use implicitly---post-application repair as a convergence
mechanism---and provides both the structural conditions under which it is
sound and a calculus for verifying those conditions.

\section*{Acknowledgments}

I thank my wife, Yuval, for her patience and support throughout this research.

\appendix
\section{Errata and corrections}
\label{app:errata}

This appendix lists every statement of version~1 that version~2 refutes,
corrects or sharpens.  For each, it gives the old statement (short), the
counterexample where there is one, and the corrected result, with the names of
the Coq theorems (all in the axiom-free gate of \texttt{coq/verify.sh}).  The
body of the paper states only the corrected results.

\begin{enumerate}[leftmargin=*,label=\textbf{E\arabic*.},ref=E\arabic*]
\item \textbf{\Cref{thm:footprint-cc1} (Footprint Commutativity).}
  \emph{Old:} under WFC and decomposable repair, disjoint footprints (H1) and
  repair locality (H2) give CC1 for an independent pair.
  \emph{Counterexamples:} with $k = 0$ the events $x := 1$ and $x := 2$
  satisfy (H1) and (H2) vacuously and violate CC1
  (\coq{base_thm_footprint_cc1_refuted}); the four-state example of
  \Cref{rem:footprint-counterexamples} satisfies (H1), (H2) and raw
  commutation and violates CC1 at state~0
  (\coq{base_thm_footprint_cc1_raw_refuted}).  The old proof's final step
  moved a repair block past an event whose footprint it intersects, which
  repair locality does not allow.
  \emph{Corrected:} \Cref{thm:footprint-cc1}: (a)~under (H1), (H2), CC1 at
  $\sigma$ iff the events commute up to~$N$ after the other event's footprint
  repairs (\coq{calc_footprint_cc1_iff}); (b)~at valid states, CC1 iff
  commutation up to~$N$ (\coq{calc_footprint_cc1_valid_iff},
  \coq{calc_footprint_cc1_valid}); (c)~at every state, repair locality,
  commutation up to~$N$ and per-event footprint absorption give CC1, without
  (H1) (\coq{calc_footprint_cc1}, \coq{calc_fp_absorb_iff_sa}); and
  \Cref{cor:cc1-under-cc2}: under CC2, CC1 iff commutation up to~$N$
  (\coq{calc_cc1_iff_commute_under_cc2}).  Each hypothesis is needed
  (\coq{calc_valid_needs_commute}, \coq{calc_valid_needs_disjoint},
  \coq{calc_valid_needs_repair_local}, \coq{calc_all_needs_commute},
  \coq{calc_all_needs_absorb}, \coq{calc_all_needs_repair_local}).

\item \textbf{Remark ``Practical import''.}
  \emph{Old:} the order-fulfillment events have disjoint footprints, so the
  footprint workflow certifies the example.
  \emph{Counterexample:} the single invariant lies in every footprint of both
  events (\coq{base_rem_practical_import_refuted}).
  \emph{Corrected:} \Cref{rem:practical-import}: the workflow has a fourth
  step (per-event footprint absorption), and the example is certified by
  \Cref{thm:footprint-cc1}(c) (\coq{calc_order_credit_repair_commute},
  \coq{calc_order_cc1}).

\item \textbf{\Cref{sec:verify-cc}, pattern~1 (compensation to a canonical
  state).}
  \emph{Old:} if $\rho_R^*(\sigma) = \sigma_\bot$ for every invalid~$\sigma$,
  CC1 and CC2 hold trivially.
  \emph{Counterexamples:} $x := 1$ and $x := 2$ on valid states violate CC1
  (\coq{base_pattern1_cc1_refuted}); an event that sends an invalid state and
  the reset state to different places violates CC2
  (\coq{base_pattern1_cc2_refuted}).
  \emph{Corrected:} CC2 for~$e$ iff $N(A_e(\sigma)) = N(A_e(\sigma_\bot))$ for
  every invalid~$\sigma$ (\coq{calc_canonical_cc2_iff}); given that for both
  events, CC1 iff commutation up to~$N$ (\coq{calc_canonical_pattern},
  instance \coq{calc_canonical_instance}).

\item \textbf{\Cref{sec:verify-cc}, pattern~2 (lattice-structured
  compensation).}
  \emph{Old:} compensation that ``computes the join with the nearest valid
  state'' satisfies CC by commutativity of the join.  This was imprecise
  (``nearest valid state'' is undefined).
  \emph{Corrected:} stated precisely (events join fixed elements, $N$ a
  closure operator, $\sigma \sqsubseteq \rho_R(\sigma) \sqsubseteq N(\sigma)$)
  with a short paper proof; not mechanized.  The federated remark in the same
  item no longer claims event-order convergence on monotone cycles, which the
  companion paper's revision shows needs a further condition.

\item \textbf{\Cref{sec:example}, ``Why Compensation Design Matters''.}
  \emph{Old:} under the naive revert $\rho(\textit{approved}, b) =
  (\textit{pending}, b)$, CC1 fails at $(\textit{pending}, 0)$.
  \emph{Counterexample:} at $(\textit{pending}, 0)$ both orders give
  $(\textit{pending}, 1)$ (\coq{naive_paper_witness_fails}).
  \emph{Corrected:} CC1 fails at $(\textit{pending}, 1)$, with
  $(\textit{pending}, 2) \neq (\textit{approved}, 2)$
  (\coq{naive_cc1_fails}); the paper's stream diverges
  (\coq{naive_stream_diverges}); the naive revert also violates CC2 at
  $(\textit{approved}, 1)$ (\coq{naive_cc2_fails}).  The closing paragraph
  no longer cites ``state-determined compensation'' as a sufficient
  condition (canonical compensation is a special case and is not sufficient,
  E3); it names the absorption and commutation facts that
  \Cref{thm:footprint-cc1} uses.  The exhaustive check covers all 12 states
  (\coq{of_cc1}, \coq{of_cc2}, \coq{of_processors}).

\item \textbf{Remark ``What Fails'' (after \Cref{thm:necessity}).}
  \emph{Old:} $R_\infty$ satisfies CC because ``there is only one event per
  set, so order independence is vacuous''.  The conclusion is true, the
  reason is not (CC1 is a property of the registry over all events, and
  $\mathrm{apply}(e_n, -)$ was defined only at~0).
  \emph{Corrected:} $\rho^*$ is constantly~0, so CC1 and CC2 hold for every
  total extension of $\mathrm{apply}(e_n, -)$ (\coq{ri_cc_any_extension},
  \coq{ri_what_fails}); stronger forms of the UBC failure:
  \coq{ri_depth_exact}, \coq{ri_no_ubc}.  \Cref{thm:necessity} itself is
  unchanged (\coq{thm_necessity}).

\item \textbf{\Cref{thm:convergence}(c) (Agreement).}
  \emph{Old:} if $E_{P_1}(t) = E_{P_2}(t)$ then $\Psi_{P_1}(t) =
  \Psi_{P_2}(t)$.
  \emph{Counterexample:} one processor has applied the single received event
  and the other has not yet (\coq{base_thm_convergence_c_counterexample}).
  \emph{Corrected:} processors that have \emph{finished processing} equal
  received sets agree, at the same or different times
  (\coq{base_thm_convergence_c}, \coq{stream_agreement}); parts (a) and~(b)
  are stated for finished processors (\coq{base_thm_convergence},
  \coq{stream_convergence}).  \Cref{def:processor} now defines
  ``finished'' and ``fair''.

\item \textbf{Closing sentence of \Cref{thm:convergence} and
  \Cref{rem:infinite-streams}.}
  \emph{Old:} with eventual delivery, disagreements are transient: any event
  causing processors to differ is eventually received by both, restoring
  agreement.
  \emph{Counterexample:} on an infinite stream, a processor one event behind
  another is finished at every time and disagrees at every time
  (\coq{base_thm_convergence_transient_counterexample}).
  \emph{Corrected:} the sentence is removed; \Cref{rem:finished} and
  \Cref{rem:infinite-streams} state what holds: (c) as corrected, and
  fixed-point agreement for finite streams.

\item \textbf{\Cref{cor:quiescent}.}
  \emph{Old:} for a finite stream the processors eventually agree forever.
  The proof assumed both processors had finished processing at~$T_N$.
  \emph{Corrected:} the same conclusion for fair processors
  (\coq{base_cor_quiescent}, \coq{base_cor_quiescent_nat}).

\item \textbf{``CC is necessary'' (abstract, contributions,
  \Cref{sec:necessity}, the remark after \Cref{prop:cc-necessary},
  conclusion).}
  \emph{Old:} CC is necessary for agreement, and each condition is
  independently necessary for its guarantee.  \Cref{prop:cc-necessary} is
  existential and remains true (\coq{prop_cc_necessary}), but CC as stated is
  not necessary in the paper's causal setting.
  \emph{Counterexamples:} under causal enabledness CC1 can fail at a
  reachable configuration while normal forms are unique
  (\coq{masked_cc1}); with non-canonical $\rho^*$ the same
  (\coq{rho_star_qualifier}).
  \emph{Corrected:} \Cref{thm:cc-exact}: under free delivery and canonical
  repair, unique normal forms from~$\sigma_0$ iff CC1 and CC2 on the states
  reachable from~$\sigma_0$ (\coq{cc_exact_from}, \coq{wfc_cc_exact_from});
  under any enabledness, confluence iff joinability at reachable
  configurations (\coq{jc_exact}); \Cref{rem:converse-scope}.

\item \textbf{\Cref{thm:strong-absorption}.}
  \emph{Old:} strong absorption implies CC2.
  \emph{Sharpened:} under WFC the two are equivalent
  (\coq{strong_absorption_iff_cc2}; forward direction
  \coq{base_thm_strong_absorption}; per event
  \coq{calc_cc2_strong_absorption}).

\item \textbf{Three-regime hierarchy (\Cref{sec:intro},
  \Cref{sec:design-space}, Related Work).}
  \emph{Old:} the regimes form a containment hierarchy of progressively
  weaker requirements, and normalization confluence subsumes invariant
  confluence.  This was imprecise: an invariant-confluent system needs no
  compensation but its operations need not commute, so it can violate CC1.
  \emph{Corrected:} normalization confluence contains the
  invariant-confluent systems that converge under reordering; the op-based
  CRDT side is exact (\coq{compensation_free_exact},
  \coq{cmrdt_governed_SEC}, \coq{witness_converges},
  \coq{witness_not_cmrdt}).

\item \textbf{Sharpened or newly mechanized, no statement refuted.}
  \Cref{def:rhostar}: $\rho_R^*$ is constructed from WFC, needs at most
  $\Phi(\sigma)$ steps, and is independent of the measure
  (\coq{base_def_rhostar}, \coq{base_def_rhostar_least_unique},
  \coq{rho_star_canonical}).  \Cref{lem:termination}: the bound is attained
  (\coq{base_lem_termination_bound}, \coq{lv_termination_bound_tight}).
  \Cref{cor:infinite} (domain independence) is mechanized at the
  rewrite-system and stream levels (\coq{governance_wf_confluent},
  \coq{causal_governance_wf_confluent}, \coq{base_cor_infinite},
  \coq{base_cor_infinite_quiescent}, \coq{zw_confluent}) and now also covers
  \Cref{cor:quiescent}; its paragraph on widening no longer claims
  uniqueness of the normal form without WFC.  \Cref{thm:product}: CC1 lifts
  as an equivalence with the componentwise conditions
  (\coq{base_thm_product}, \coq{base_thm_product_decomposable}).  The
  Applications paragraphs no longer assert CC for reset and per-constraint
  remediation; they point to the corrected checks.  The remark
  ``Machine-checked'' (\Cref{sec:mechanization}) lists the Coq theorem
  names; version~1's phrase ``footprint disjointness implies
  order-independent commutation'' referred to state-variable footprints
  (\coq{disjoint_events_commute}), not to \Cref{thm:footprint-cc1}.
\end{enumerate}

\begin{thebibliography}{99}

\bibitem{newman1942}
M.~H.~A. Newman,
``On theories with a combinatorial definition of `equivalence',''
\emph{Annals of Mathematics}, vol.~43, no.~2, pp.~223--243, 1942.

\bibitem{shapiro2011crdt}
M.~Shapiro, N.~Pregui\c{c}a, C.~Baquero, and M.~Zawirski,
``A comprehensive study of convergent and commutative replicated data types,''
\emph{INRIA Research Report RR-7506}, 2011.

\bibitem{hellerstein2010declarative}
J.~M.~Hellerstein,
``The declarative imperative: experiences and conjectures in distributed logic,''
\emph{SIGMOD Record}, vol.~39, no.~1, pp.~5--19, 2010.

\bibitem{ameloot2013relational}
T.~J.~Ameloot, F.~Neven, and J.~Van~den~Bussche,
``Relational transducers for declarative networking,''
\emph{Journal of the ACM}, vol.~60, no.~2, pp.~1--38, 2013.

\bibitem{bailis2014coordination}
P.~Bailis, A.~Fekete, M.~J.~Franklin, A.~Ghodsi, J.~M.~Hellerstein, and I.~Stoica,
``Coordination avoidance in database systems,''
\emph{Proc.\ VLDB Endowment}, vol.~8, no.~3, pp.~185--196, 2014.

\bibitem{baader1998term}
F.~Baader and T.~Nipkow,
\emph{Term Rewriting and All That},
Cambridge University Press, 1998.

\bibitem{vogels2009eventually}
W.~Vogels,
``Eventually consistent,''
\emph{Communications of the ACM}, vol.~52, no.~1, pp.~40--44, 2009.

\bibitem{garcia1987sagas}
H.~Garcia-Molina and K.~Salem,
``Sagas,''
\emph{Proc.\ ACM SIGMOD International Conference on Management of Data}, pp.~249--259, 1987.

\bibitem{lee1987synchronous}
E.~A.~Lee and D.~G.~Messerschmitt,
``Synchronous data flow,''
\emph{Proceedings of the IEEE}, vol.~75, no.~9, pp.~1235--1245, 1987.

\bibitem{kahn1974semantics}
G.~Kahn,
``The semantics of a simple language for parallel programming,''
\emph{Information Processing}, pp.~471--475, 1974.

\bibitem{akidau2015dataflow}
T.~Akidau, R.~Bradshaw, C.~Chambers, et~al.,
``The Dataflow Model: A practical approach to balancing correctness, latency, and cost in massive-scale, unbounded, out-of-order data processing,''
\emph{Proc.\ VLDB Endowment}, vol.~8, no.~12, pp.~1792--1803, 2015.

\bibitem{lamport1978time}
L.~Lamport,
``Time, clocks, and the ordering of events in a distributed system,''
\emph{Communications of the ACM}, vol.~21, no.~7, pp.~558--565, 1978.

\bibitem{blackwell2026federated}
Dayna~Blackwell,
``Normalization Confluence in Federated Registry Networks,''
Zenodo, 2026.
\textsc{doi}: \href{https://doi.org/10.5281/zenodo.18677400}{10.5281/zenodo.18677400}.

\bibitem{cousot1977}
P.~Cousot and R.~Cousot,
``Abstract interpretation: a unified lattice model for static analysis of programs by construction or approximation of fixpoints,''
\emph{Proc.\ 4th ACM SIGACT-SIGPLAN Symp.\ on Principles of Programming Languages (POPL)}, pp.~238--252, 1977.

\bibitem{balegas2015indigo}
V.~Balegas, S.~Duarte, C.~Ferreira, R.~Rodrigues, N.~Pregui\c{c}a, M.~Najafzadeh, and M.~Shapiro,
``Putting consistency back into eventual consistency,''
\emph{Proc.\ 10th European Conf.\ on Computer Systems (EuroSys)}, 2015.

\bibitem{balegas2019ipa}
V.~Balegas, S.~Duarte, C.~Ferreira, R.~Rodrigues, and N.~Pregui\c{c}a,
``IPA: Invariant-preserving applications for weakly consistent replicated databases,''
\emph{Proc.\ VLDB Endowment}, vol.~12, no.~4, 2018.

\bibitem{meseguer1992}
J.~Meseguer,
``Conditional rewriting logic as a unified model of concurrency,''
\emph{Theoretical Computer Science}, vol.~96, no.~1, pp.~73--155, 1992.

\bibitem{duran2012coherence}
F.~Dur\'an and J.~Meseguer,
``On the Church-Rosser and coherence properties of conditional order-sorted rewrite theories,''
\emph{Journal of Logic and Algebraic Programming}, vol.~81, no.~7--8, pp.~816--850, 2012.

\bibitem{dijkstra1974}
E.~W.~Dijkstra,
``Self-stabilizing systems in spite of distributed control,''
\emph{Communications of the ACM}, vol.~17, no.~11, pp.~643--644, 1974.

\bibitem{gomes2017verifying}
V.~B.~F.~Gomes, M.~Kleppmann, D.~P.~Mulligan, and A.~R.~Beresford,
``Verifying strong eventual consistency in distributed systems,''
\emph{Proc.\ ACM on Programming Languages}, vol.~1, no.~OOPSLA, art.~109, 2017.

\end{thebibliography}

\end{document}
